\begin{table}[H]
\centering
\caption{Row 13 Results (No pre-training before pruning). \ifshowcaption{}Caption: ``High-flying tax lawter Cynthia Gibson Beerbower, 65, was found by her housekeeper at her apartment in Fulham, west London, in July this year . A high-flying tax lawyer who worked in the White House was found dead in the bath of her London home after taking an overdose while going through a 'painful' divorce. American-born Cynthia Gibson Beerbower, 65, was found by her housekeeper at her multi-million-pound apartment in Fulham, west London, on July 26 this year. The 'highly intelligent' former Wall Street lawyer had left three notes in her spare bedroom alongside a piece of card which read 'I can't do this anymore', an inquest has heard. The hearing was also told how the note included a request for Mrs Beerbower be buried in Cambridge, where she attended university. Mrs Beerbower had apparently been offered several ambassadorial roles by President Barack Obama in the weeks before her death and was planning a future with her new fiancé. But the hearing at Westminster Coroner's Court heard that Mrs Beerbower was struggling to cope with an ongoing divorce battle and feared her family 'wished her dead'. Her estranged husband also described how Mrs Beerbower had a history of mental illness and had threatened suicide several times. During the hearing, Coroner Fiona Wilcox told the court that a note left by the lawyer said: 'I can't do this any more, I hate you, I am sick'. The court also heard evidence from medical professionals who said Mrs Beerbower had started taking medication that she had not been prescribed. During the hearing, Dr Marios Pierides, a psychiatrist at the Nightingale Hospital in Marylebone, said Mrs Beerbower had visited him before her death. He described how she felt 'lonely' and 'isolated' while going through the process of divorcing her husband John Beerbower. He told the court: 'Cynthia was on various drugs at different time. 'When she came to me she was being treated with various medication, she was on anti-depressants and sleeping pills. 'She explained to me she was going through a very painful period of divorce. Also she was feeling quite lonely and isolated in her residence in Fulham. Dr Pierides added: 'She would say the family wished her dead. But she was, from my point of view, a life affirming person who enjoyed things until the awful circumstances around the divorce happened.' He continued: 'At the same time she was angry and lonely at her circumstances.' The court heard Mrs Beerbower had started taking Phentermine, an appetite suppressor, as well as Diazepam. Neither had been prescribed by the psychiatrist. Dr Pierides told the hearing that the drugs have the capability to 'alter' a person's mental state, adding: 'It could promote psychosis or mood swings. 'It can alter your mental state and you might not know what you're doing.' The American-born Cambridge graduate (left) had been offered several US ambassadorial roles by Barack Obama in the weeks before her death and was planning a future with her new partner, David Shaw (right) Mrs Beerbower's estranged husband, whom she married in 1971, told the court that Mrs Beerbower had threatened suicide in the past. He also said she had previously 'over-medicated' and had mental health problems dating back to the 1970s. Mr Beebower told the inquest: 'I think there were some issues that go all the way back to the middle of the 70s. It was between two and three decades that she had been seeing people. ' He said that she became suicidal on occasions, adding: 'She would become very upset or angry and would make these kinds of statements and, for example, lock herself in the bathroom. 'As far as I'm aware she never undertook any action on any of these threats.' The inquest heard Mrs Beerbower was struggling to cope with her divorce from husband John and feared her family 'wished her dead'. Pictured: Mr Beerbower and the couple's daughter Sarah outside Westminster Coroner's Court . Despite her divorce not being finalised, Mrs Beerbower become engaged to David Shaw in May this year, the inquest heard. Mr Shaw described her as a 'highly intelligent women' and said they had started to make plans for the future. He said: 'We were planning for the next 20 years. We were planning about what countries we wanted to go, because she was being offered ambassadorial roles by President Obama.' Mrs Beerbower - one of the few female partners in a Wall Street law firm - had joined the first Clinton administration in 1993. She served as International Tax Counsel under Secretary Lloyd Bentsen before becoming Deputy Assistant Secretary for Tax Policy under Secretary Robert Rubin. She left US Government service in late 1996. Mrs Beerbower retired from the financial business in 2005 and moved to France with her daughter Sarah. Mrs Beerbower was found dead by her housekeeper at her apartment in Fulham, west London (pictured) She and her husband acquired an olive farm in Chateauneuf de Grasse, in the south of France in 2006. From 2012 onward, Cynthia divided her time between homes in Toronto, New York City, the South of France and London, the inquest heard. Home Office pathologist Dr Jon Van Der Walt gave the cause of death as Oxidon, Phetermine and paracetemol toxicity. He said the drugs could have been enough to create a loss of consciousness, or fatal arrhythmia. The coroner said she could not give a verdict of suicide because the drugs Mrs Beerbower was taking could have affected her decision-making. Dr Wilcox said: 'At the time of death she was under the influence of Phentermine - it may well have affected her ability to have clear intentions, or even think clearly in any sensible way.' She instead recorded a narrative verdict, adding: 'She died as a result of poly-drug overdose while under the influence of Phentermine.' For confidential support call the Samaritans in the UK on 08457 90 90 90, visit a local Samaritans branch or click here for details.''.\fi}
\label{tab:evaluation-pipeline-13-untrained}
\begin{tabular}{lllrlllll}
\toprule
 &  & Perplexity & ROGUE-L F1 Score (%) & TTFT (ms) & Run Time (ms) & Output Tokens & Time per Output Token (ms) & Generated Text \\
Trained & Pruning Percentage &  &  &  &  &  &  &  \\
\midrule
nan & 0% & \bfseries 23.1183 & \bfseries 100.0000\% & 12.846 & 493.577 & 1,283 & 0.385 & High-flying tax lawter Cynthia Gibson Beerbower, 65, was found by her housekeeper at her apartment in Fulham, west London, in July this year . A high-flying tax lawyer who worked in the White House was found dead in the bath of her London home after taking an overdose while going through a 'painful' divorce. American-born Cynthia Gibson Beerbower, 65, was found by her housekeeper at her multi-million-pound apartment in Fulham, west London, on July 26 this year. The 'highly intelligent' former Wall Street lawyer had left three notes in her spare bedroom alongside a piece of card which read 'I can't do this anymore', an inquest has heard. The hearing was also told how the note included a request for Mrs Beerbower be buried in Cambridge, where she attended university. Mrs Beerbower had apparently been offered several ambassadorial roles by President Barack Obama in the weeks before her death and was planning a future with her new fiancé. But the hearing at Westminster Coroner's Court heard that Mrs Beerbower was struggling to cope with an ongoing divorce battle and feared her family 'wished her dead'. Her estranged husband also described how Mrs Beerbower had a history of mental illness and had threatened suicide several times. During the hearing, Coroner Fiona Wilcox told the court that a note left by the lawyer said: 'I can't do this any more, I hate you, I am sick'. The court also heard evidence from medical professionals who said Mrs Beerbower had started taking medication that she had not been prescribed. During the hearing, Dr Marios Pierides, a psychiatrist at the Nightingale Hospital in Marylebone, said Mrs Beerbower had visited him before her death. He described how she felt 'lonely' and 'isolated' while going through the process of divorcing her husband John Beerbower. He told the court: 'Cynthia was on various drugs at different time. 'When she came to me she was being treated with various medication, she was on anti-depressants and sleeping pills. 'She explained to me she was going through a very painful period of divorce. Also she was feeling quite lonely and isolated in her residence in Fulham. Dr Pierides added: 'She would say the family wished her dead. But she was, from my point of view, a life affirming person who enjoyed things until the awful circumstances around the divorce happened.' He continued: 'At the same time she was angry and lonely at her circumstances.' The court heard Mrs Beerbower had started taking Phentermine, an appetite suppressor, as well as Diazepam. Neither had been prescribed by the psychiatrist. Dr Pierides told the hearing that the drugs have the capability to 'alter' a person's mental state, adding: 'It could promote psychosis or mood swings. 'It can alter your mental state and you might not know what you're doing.' The American-born Cambridge graduate (left) had been offered several US ambassadorial roles by Barack Obama in the weeks before her death and was planning a future with her new partner, David Shaw (right) Mrs Beerbower's estranged husband, whom she married in 1971, told the court that Mrs Beerbower had threatened suicide in the past. He also said she had previously 'over-medicated' and had mental health problems dating back to the 1970s. Mr Beebower told the inquest: 'I think there were some issues that go all the way back to the middle of the 70s. It was between two and three decades that she had been seeing people. ' He said that she became suicidal on occasions, adding: 'She would become very upset or angry and would make these kinds of statements and, for example, lock herself in the bathroom. 'As far as I'm aware she never undertook any action on any of these threats.' The inquest heard Mrs Beerbower was struggling to cope with her divorce from husband John and feared her family 'wished her dead'. Pictured: Mr Beerbower and the couple's daughter Sarah outside Westminster Coroner's Court . Despite her divorce not being finalised, Mrs Beerbower become engaged to David Shaw in May this year, the inquest heard. Mr Shaw described her as a 'highly intelligent women' and said they had started to make plans for the future. He said: 'We were planning for the next 20 years. We were planning about what countries we wanted to go, because she was being offered ambassadorial roles by President Obama.' Mrs Beerbower - one of the few female partners in a Wall Street law firm - had joined the first Clinton administration in 1993. She served as International Tax Counsel under Secretary Lloyd Bentsen before becoming Deputy Assistant Secretary for Tax Policy under Secretary Robert Rubin. She left US Government service in late 1996. Mrs Beerbower retired from the financial business in 2005 and moved to France with her daughter Sarah. Mrs Beerbower was found dead by her housekeeper at her apartment in Fulham, west London (pictured) She and her husband acquired an olive farm in Chateauneuf de Grasse, in the south of France in 2006. From 2012 onward, Cynthia divided her time between homes in Toronto, New York City, the South of France and London, the inquest heard. Home Office pathologist Dr Jon Van Der Walt gave the cause of death as Oxidon, Phetermine and paracetemol toxicity. He said the drugs could have been enough to create a loss of consciousness, or fatal arrhythmia. The coroner said she could not give a verdict of suicide because the drugs Mrs Beerbower was taking could have affected her decision-making. Dr Wilcox said: 'At the time of death she was under the influence of Phentermine - it may well have affected her ability to have clear intentions, or even think clearly in any sensible way.' She instead recorded a narrative verdict, adding: 'She died as a result of poly-drug overdose while under the influence of Phentermine.' For confidential support call the Samaritans in the UK on 08457 90 90 90, visit a local Samaritans branch or click here for details.
 \\
\cline{1-9}
False & 2.5% & \color{red} 1,202,604.2842 & \bfseries \color{red} 100.0000\% & 12.106 & 475.568 & 1,283 & 0.371 & High-flying tax lawter Cynthia Gibson Beerbower, 65, was found by her housekeeper at her apartment in Fulham, west London, in July this year . A high-flying tax lawyer who worked in the White House was found dead in the bath of her London home after taking an overdose while going through a 'painful' divorce. American-born Cynthia Gibson Beerbower, 65, was found by her housekeeper at her multi-million-pound apartment in Fulham, west London, on July 26 this year. The 'highly intelligent' former Wall Street lawyer had left three notes in her spare bedroom alongside a piece of card which read 'I can't do this anymore', an inquest has heard. The hearing was also told how the note included a request for Mrs Beerbower be buried in Cambridge, where she attended university. Mrs Beerbower had apparently been offered several ambassadorial roles by President Barack Obama in the weeks before her death and was planning a future with her new fiancé. But the hearing at Westminster Coroner's Court heard that Mrs Beerbower was struggling to cope with an ongoing divorce battle and feared her family 'wished her dead'. Her estranged husband also described how Mrs Beerbower had a history of mental illness and had threatened suicide several times. During the hearing, Coroner Fiona Wilcox told the court that a note left by the lawyer said: 'I can't do this any more, I hate you, I am sick'. The court also heard evidence from medical professionals who said Mrs Beerbower had started taking medication that she had not been prescribed. During the hearing, Dr Marios Pierides, a psychiatrist at the Nightingale Hospital in Marylebone, said Mrs Beerbower had visited him before her death. He described how she felt 'lonely' and 'isolated' while going through the process of divorcing her husband John Beerbower. He told the court: 'Cynthia was on various drugs at different time. 'When she came to me she was being treated with various medication, she was on anti-depressants and sleeping pills. 'She explained to me she was going through a very painful period of divorce. Also she was feeling quite lonely and isolated in her residence in Fulham. Dr Pierides added: 'She would say the family wished her dead. But she was, from my point of view, a life affirming person who enjoyed things until the awful circumstances around the divorce happened.' He continued: 'At the same time she was angry and lonely at her circumstances.' The court heard Mrs Beerbower had started taking Phentermine, an appetite suppressor, as well as Diazepam. Neither had been prescribed by the psychiatrist. Dr Pierides told the hearing that the drugs have the capability to 'alter' a person's mental state, adding: 'It could promote psychosis or mood swings. 'It can alter your mental state and you might not know what you're doing.' The American-born Cambridge graduate (left) had been offered several US ambassadorial roles by Barack Obama in the weeks before her death and was planning a future with her new partner, David Shaw (right) Mrs Beerbower's estranged husband, whom she married in 1971, told the court that Mrs Beerbower had threatened suicide in the past. He also said she had previously 'over-medicated' and had mental health problems dating back to the 1970s. Mr Beebower told the inquest: 'I think there were some issues that go all the way back to the middle of the 70s. It was between two and three decades that she had been seeing people. ' He said that she became suicidal on occasions, adding: 'She would become very upset or angry and would make these kinds of statements and, for example, lock herself in the bathroom. 'As far as I'm aware she never undertook any action on any of these threats.' The inquest heard Mrs Beerbower was struggling to cope with her divorce from husband John and feared her family 'wished her dead'. Pictured: Mr Beerbower and the couple's daughter Sarah outside Westminster Coroner's Court . Despite her divorce not being finalised, Mrs Beerbower become engaged to David Shaw in May this year, the inquest heard. Mr Shaw described her as a 'highly intelligent women' and said they had started to make plans for the future. He said: 'We were planning for the next 20 years. We were planning about what countries we wanted to go, because she was being offered ambassadorial roles by President Obama.' Mrs Beerbower - one of the few female partners in a Wall Street law firm - had joined the first Clinton administration in 1993. She served as International Tax Counsel under Secretary Lloyd Bentsen before becoming Deputy Assistant Secretary for Tax Policy under Secretary Robert Rubin. She left US Government service in late 1996. Mrs Beerbower retired from the financial business in 2005 and moved to France with her daughter Sarah. Mrs Beerbower was found dead by her housekeeper at her apartment in Fulham, west London (pictured) She and her husband acquired an olive farm in Chateauneuf de Grasse, in the south of France in 2006. From 2012 onward, Cynthia divided her time between homes in Toronto, New York City, the South of France and London, the inquest heard. Home Office pathologist Dr Jon Van Der Walt gave the cause of death as Oxidon, Phetermine and paracetemol toxicity. He said the drugs could have been enough to create a loss of consciousness, or fatal arrhythmia. The coroner said she could not give a verdict of suicide because the drugs Mrs Beerbower was taking could have affected her decision-making. Dr Wilcox said: 'At the time of death she was under the influence of Phentermine - it may well have affected her ability to have clear intentions, or even think clearly in any sensible way.' She instead recorded a narrative verdict, adding: 'She died as a result of poly-drug overdose while under the influence of Phentermine.' For confidential support call the Samaritans in the UK on 08457 90 90 90, visit a local Samaritans branch or click here for details.Ồ \\
\cline{1-9}
False & 5.0% & 1,862,629.9526 & \bfseries \color{red} 100.0000\% & 12.223 & \bfseries \color{red} 473.186 & 1,283 & \bfseries \color{red} 0.369 & High-flying tax lawter Cynthia Gibson Beerbower, 65, was found by her housekeeper at her apartment in Fulham, west London, in July this year . A high-flying tax lawyer who worked in the White House was found dead in the bath of her London home after taking an overdose while going through a 'painful' divorce. American-born Cynthia Gibson Beerbower, 65, was found by her housekeeper at her multi-million-pound apartment in Fulham, west London, on July 26 this year. The 'highly intelligent' former Wall Street lawyer had left three notes in her spare bedroom alongside a piece of card which read 'I can't do this anymore', an inquest has heard. The hearing was also told how the note included a request for Mrs Beerbower be buried in Cambridge, where she attended university. Mrs Beerbower had apparently been offered several ambassadorial roles by President Barack Obama in the weeks before her death and was planning a future with her new fiancé. But the hearing at Westminster Coroner's Court heard that Mrs Beerbower was struggling to cope with an ongoing divorce battle and feared her family 'wished her dead'. Her estranged husband also described how Mrs Beerbower had a history of mental illness and had threatened suicide several times. During the hearing, Coroner Fiona Wilcox told the court that a note left by the lawyer said: 'I can't do this any more, I hate you, I am sick'. The court also heard evidence from medical professionals who said Mrs Beerbower had started taking medication that she had not been prescribed. During the hearing, Dr Marios Pierides, a psychiatrist at the Nightingale Hospital in Marylebone, said Mrs Beerbower had visited him before her death. He described how she felt 'lonely' and 'isolated' while going through the process of divorcing her husband John Beerbower. He told the court: 'Cynthia was on various drugs at different time. 'When she came to me she was being treated with various medication, she was on anti-depressants and sleeping pills. 'She explained to me she was going through a very painful period of divorce. Also she was feeling quite lonely and isolated in her residence in Fulham. Dr Pierides added: 'She would say the family wished her dead. But she was, from my point of view, a life affirming person who enjoyed things until the awful circumstances around the divorce happened.' He continued: 'At the same time she was angry and lonely at her circumstances.' The court heard Mrs Beerbower had started taking Phentermine, an appetite suppressor, as well as Diazepam. Neither had been prescribed by the psychiatrist. Dr Pierides told the hearing that the drugs have the capability to 'alter' a person's mental state, adding: 'It could promote psychosis or mood swings. 'It can alter your mental state and you might not know what you're doing.' The American-born Cambridge graduate (left) had been offered several US ambassadorial roles by Barack Obama in the weeks before her death and was planning a future with her new partner, David Shaw (right) Mrs Beerbower's estranged husband, whom she married in 1971, told the court that Mrs Beerbower had threatened suicide in the past. He also said she had previously 'over-medicated' and had mental health problems dating back to the 1970s. Mr Beebower told the inquest: 'I think there were some issues that go all the way back to the middle of the 70s. It was between two and three decades that she had been seeing people. ' He said that she became suicidal on occasions, adding: 'She would become very upset or angry and would make these kinds of statements and, for example, lock herself in the bathroom. 'As far as I'm aware she never undertook any action on any of these threats.' The inquest heard Mrs Beerbower was struggling to cope with her divorce from husband John and feared her family 'wished her dead'. Pictured: Mr Beerbower and the couple's daughter Sarah outside Westminster Coroner's Court . Despite her divorce not being finalised, Mrs Beerbower become engaged to David Shaw in May this year, the inquest heard. Mr Shaw described her as a 'highly intelligent women' and said they had started to make plans for the future. He said: 'We were planning for the next 20 years. We were planning about what countries we wanted to go, because she was being offered ambassadorial roles by President Obama.' Mrs Beerbower - one of the few female partners in a Wall Street law firm - had joined the first Clinton administration in 1993. She served as International Tax Counsel under Secretary Lloyd Bentsen before becoming Deputy Assistant Secretary for Tax Policy under Secretary Robert Rubin. She left US Government service in late 1996. Mrs Beerbower retired from the financial business in 2005 and moved to France with her daughter Sarah. Mrs Beerbower was found dead by her housekeeper at her apartment in Fulham, west London (pictured) She and her husband acquired an olive farm in Chateauneuf de Grasse, in the south of France in 2006. From 2012 onward, Cynthia divided her time between homes in Toronto, New York City, the South of France and London, the inquest heard. Home Office pathologist Dr Jon Van Der Walt gave the cause of death as Oxidon, Phetermine and paracetemol toxicity. He said the drugs could have been enough to create a loss of consciousness, or fatal arrhythmia. The coroner said she could not give a verdict of suicide because the drugs Mrs Beerbower was taking could have affected her decision-making. Dr Wilcox said: 'At the time of death she was under the influence of Phentermine - it may well have affected her ability to have clear intentions, or even think clearly in any sensible way.' She instead recorded a narrative verdict, adding: 'She died as a result of poly-drug overdose while under the influence of Phentermine.' For confidential support call the Samaritans in the UK on 08457 90 90 90, visit a local Samaritans branch or click here for details.targetContext \\
\cline{1-9}
False & 7.5% & 1,862,629.9526 & \bfseries \color{red} 100.0000\% & 12.122 & 473.320 & 1,283 & 0.369 & High-flying tax lawter Cynthia Gibson Beerbower, 65, was found by her housekeeper at her apartment in Fulham, west London, in July this year . A high-flying tax lawyer who worked in the White House was found dead in the bath of her London home after taking an overdose while going through a 'painful' divorce. American-born Cynthia Gibson Beerbower, 65, was found by her housekeeper at her multi-million-pound apartment in Fulham, west London, on July 26 this year. The 'highly intelligent' former Wall Street lawyer had left three notes in her spare bedroom alongside a piece of card which read 'I can't do this anymore', an inquest has heard. The hearing was also told how the note included a request for Mrs Beerbower be buried in Cambridge, where she attended university. Mrs Beerbower had apparently been offered several ambassadorial roles by President Barack Obama in the weeks before her death and was planning a future with her new fiancé. But the hearing at Westminster Coroner's Court heard that Mrs Beerbower was struggling to cope with an ongoing divorce battle and feared her family 'wished her dead'. Her estranged husband also described how Mrs Beerbower had a history of mental illness and had threatened suicide several times. During the hearing, Coroner Fiona Wilcox told the court that a note left by the lawyer said: 'I can't do this any more, I hate you, I am sick'. The court also heard evidence from medical professionals who said Mrs Beerbower had started taking medication that she had not been prescribed. During the hearing, Dr Marios Pierides, a psychiatrist at the Nightingale Hospital in Marylebone, said Mrs Beerbower had visited him before her death. He described how she felt 'lonely' and 'isolated' while going through the process of divorcing her husband John Beerbower. He told the court: 'Cynthia was on various drugs at different time. 'When she came to me she was being treated with various medication, she was on anti-depressants and sleeping pills. 'She explained to me she was going through a very painful period of divorce. Also she was feeling quite lonely and isolated in her residence in Fulham. Dr Pierides added: 'She would say the family wished her dead. But she was, from my point of view, a life affirming person who enjoyed things until the awful circumstances around the divorce happened.' He continued: 'At the same time she was angry and lonely at her circumstances.' The court heard Mrs Beerbower had started taking Phentermine, an appetite suppressor, as well as Diazepam. Neither had been prescribed by the psychiatrist. Dr Pierides told the hearing that the drugs have the capability to 'alter' a person's mental state, adding: 'It could promote psychosis or mood swings. 'It can alter your mental state and you might not know what you're doing.' The American-born Cambridge graduate (left) had been offered several US ambassadorial roles by Barack Obama in the weeks before her death and was planning a future with her new partner, David Shaw (right) Mrs Beerbower's estranged husband, whom she married in 1971, told the court that Mrs Beerbower had threatened suicide in the past. He also said she had previously 'over-medicated' and had mental health problems dating back to the 1970s. Mr Beebower told the inquest: 'I think there were some issues that go all the way back to the middle of the 70s. It was between two and three decades that she had been seeing people. ' He said that she became suicidal on occasions, adding: 'She would become very upset or angry and would make these kinds of statements and, for example, lock herself in the bathroom. 'As far as I'm aware she never undertook any action on any of these threats.' The inquest heard Mrs Beerbower was struggling to cope with her divorce from husband John and feared her family 'wished her dead'. Pictured: Mr Beerbower and the couple's daughter Sarah outside Westminster Coroner's Court . Despite her divorce not being finalised, Mrs Beerbower become engaged to David Shaw in May this year, the inquest heard. Mr Shaw described her as a 'highly intelligent women' and said they had started to make plans for the future. He said: 'We were planning for the next 20 years. We were planning about what countries we wanted to go, because she was being offered ambassadorial roles by President Obama.' Mrs Beerbower - one of the few female partners in a Wall Street law firm - had joined the first Clinton administration in 1993. She served as International Tax Counsel under Secretary Lloyd Bentsen before becoming Deputy Assistant Secretary for Tax Policy under Secretary Robert Rubin. She left US Government service in late 1996. Mrs Beerbower retired from the financial business in 2005 and moved to France with her daughter Sarah. Mrs Beerbower was found dead by her housekeeper at her apartment in Fulham, west London (pictured) She and her husband acquired an olive farm in Chateauneuf de Grasse, in the south of France in 2006. From 2012 onward, Cynthia divided her time between homes in Toronto, New York City, the South of France and London, the inquest heard. Home Office pathologist Dr Jon Van Der Walt gave the cause of death as Oxidon, Phetermine and paracetemol toxicity. He said the drugs could have been enough to create a loss of consciousness, or fatal arrhythmia. The coroner said she could not give a verdict of suicide because the drugs Mrs Beerbower was taking could have affected her decision-making. Dr Wilcox said: 'At the time of death she was under the influence of Phentermine - it may well have affected her ability to have clear intentions, or even think clearly in any sensible way.' She instead recorded a narrative verdict, adding: 'She died as a result of poly-drug overdose while under the influence of Phentermine.' For confidential support call the Samaritans in the UK on 08457 90 90 90, visit a local Samaritans branch or click here for details.Vous \\
\cline{1-9}
False & 10.0% & 2,391,664.2010 & \bfseries \color{red} 100.0000\% & 12.055 & 479.872 & 1,283 & 0.374 & High-flying tax lawter Cynthia Gibson Beerbower, 65, was found by her housekeeper at her apartment in Fulham, west London, in July this year . A high-flying tax lawyer who worked in the White House was found dead in the bath of her London home after taking an overdose while going through a 'painful' divorce. American-born Cynthia Gibson Beerbower, 65, was found by her housekeeper at her multi-million-pound apartment in Fulham, west London, on July 26 this year. The 'highly intelligent' former Wall Street lawyer had left three notes in her spare bedroom alongside a piece of card which read 'I can't do this anymore', an inquest has heard. The hearing was also told how the note included a request for Mrs Beerbower be buried in Cambridge, where she attended university. Mrs Beerbower had apparently been offered several ambassadorial roles by President Barack Obama in the weeks before her death and was planning a future with her new fiancé. But the hearing at Westminster Coroner's Court heard that Mrs Beerbower was struggling to cope with an ongoing divorce battle and feared her family 'wished her dead'. Her estranged husband also described how Mrs Beerbower had a history of mental illness and had threatened suicide several times. During the hearing, Coroner Fiona Wilcox told the court that a note left by the lawyer said: 'I can't do this any more, I hate you, I am sick'. The court also heard evidence from medical professionals who said Mrs Beerbower had started taking medication that she had not been prescribed. During the hearing, Dr Marios Pierides, a psychiatrist at the Nightingale Hospital in Marylebone, said Mrs Beerbower had visited him before her death. He described how she felt 'lonely' and 'isolated' while going through the process of divorcing her husband John Beerbower. He told the court: 'Cynthia was on various drugs at different time. 'When she came to me she was being treated with various medication, she was on anti-depressants and sleeping pills. 'She explained to me she was going through a very painful period of divorce. Also she was feeling quite lonely and isolated in her residence in Fulham. Dr Pierides added: 'She would say the family wished her dead. But she was, from my point of view, a life affirming person who enjoyed things until the awful circumstances around the divorce happened.' He continued: 'At the same time she was angry and lonely at her circumstances.' The court heard Mrs Beerbower had started taking Phentermine, an appetite suppressor, as well as Diazepam. Neither had been prescribed by the psychiatrist. Dr Pierides told the hearing that the drugs have the capability to 'alter' a person's mental state, adding: 'It could promote psychosis or mood swings. 'It can alter your mental state and you might not know what you're doing.' The American-born Cambridge graduate (left) had been offered several US ambassadorial roles by Barack Obama in the weeks before her death and was planning a future with her new partner, David Shaw (right) Mrs Beerbower's estranged husband, whom she married in 1971, told the court that Mrs Beerbower had threatened suicide in the past. He also said she had previously 'over-medicated' and had mental health problems dating back to the 1970s. Mr Beebower told the inquest: 'I think there were some issues that go all the way back to the middle of the 70s. It was between two and three decades that she had been seeing people. ' He said that she became suicidal on occasions, adding: 'She would become very upset or angry and would make these kinds of statements and, for example, lock herself in the bathroom. 'As far as I'm aware she never undertook any action on any of these threats.' The inquest heard Mrs Beerbower was struggling to cope with her divorce from husband John and feared her family 'wished her dead'. Pictured: Mr Beerbower and the couple's daughter Sarah outside Westminster Coroner's Court . Despite her divorce not being finalised, Mrs Beerbower become engaged to David Shaw in May this year, the inquest heard. Mr Shaw described her as a 'highly intelligent women' and said they had started to make plans for the future. He said: 'We were planning for the next 20 years. We were planning about what countries we wanted to go, because she was being offered ambassadorial roles by President Obama.' Mrs Beerbower - one of the few female partners in a Wall Street law firm - had joined the first Clinton administration in 1993. She served as International Tax Counsel under Secretary Lloyd Bentsen before becoming Deputy Assistant Secretary for Tax Policy under Secretary Robert Rubin. She left US Government service in late 1996. Mrs Beerbower retired from the financial business in 2005 and moved to France with her daughter Sarah. Mrs Beerbower was found dead by her housekeeper at her apartment in Fulham, west London (pictured) She and her husband acquired an olive farm in Chateauneuf de Grasse, in the south of France in 2006. From 2012 onward, Cynthia divided her time between homes in Toronto, New York City, the South of France and London, the inquest heard. Home Office pathologist Dr Jon Van Der Walt gave the cause of death as Oxidon, Phetermine and paracetemol toxicity. He said the drugs could have been enough to create a loss of consciousness, or fatal arrhythmia. The coroner said she could not give a verdict of suicide because the drugs Mrs Beerbower was taking could have affected her decision-making. Dr Wilcox said: 'At the time of death she was under the influence of Phentermine - it may well have affected her ability to have clear intentions, or even think clearly in any sensible way.' She instead recorded a narrative verdict, adding: 'She died as a result of poly-drug overdose while under the influence of Phentermine.' For confidential support call the Samaritans in the UK on 08457 90 90 90, visit a local Samaritans branch or click here for details. incorporación \\
\cline{1-9}
False & 12.5% & 2,710,110.5896 & \bfseries \color{red} 100.0000\% & 12.095 & 475.371 & 1,283 & 0.371 & High-flying tax lawter Cynthia Gibson Beerbower, 65, was found by her housekeeper at her apartment in Fulham, west London, in July this year . A high-flying tax lawyer who worked in the White House was found dead in the bath of her London home after taking an overdose while going through a 'painful' divorce. American-born Cynthia Gibson Beerbower, 65, was found by her housekeeper at her multi-million-pound apartment in Fulham, west London, on July 26 this year. The 'highly intelligent' former Wall Street lawyer had left three notes in her spare bedroom alongside a piece of card which read 'I can't do this anymore', an inquest has heard. The hearing was also told how the note included a request for Mrs Beerbower be buried in Cambridge, where she attended university. Mrs Beerbower had apparently been offered several ambassadorial roles by President Barack Obama in the weeks before her death and was planning a future with her new fiancé. But the hearing at Westminster Coroner's Court heard that Mrs Beerbower was struggling to cope with an ongoing divorce battle and feared her family 'wished her dead'. Her estranged husband also described how Mrs Beerbower had a history of mental illness and had threatened suicide several times. During the hearing, Coroner Fiona Wilcox told the court that a note left by the lawyer said: 'I can't do this any more, I hate you, I am sick'. The court also heard evidence from medical professionals who said Mrs Beerbower had started taking medication that she had not been prescribed. During the hearing, Dr Marios Pierides, a psychiatrist at the Nightingale Hospital in Marylebone, said Mrs Beerbower had visited him before her death. He described how she felt 'lonely' and 'isolated' while going through the process of divorcing her husband John Beerbower. He told the court: 'Cynthia was on various drugs at different time. 'When she came to me she was being treated with various medication, she was on anti-depressants and sleeping pills. 'She explained to me she was going through a very painful period of divorce. Also she was feeling quite lonely and isolated in her residence in Fulham. Dr Pierides added: 'She would say the family wished her dead. But she was, from my point of view, a life affirming person who enjoyed things until the awful circumstances around the divorce happened.' He continued: 'At the same time she was angry and lonely at her circumstances.' The court heard Mrs Beerbower had started taking Phentermine, an appetite suppressor, as well as Diazepam. Neither had been prescribed by the psychiatrist. Dr Pierides told the hearing that the drugs have the capability to 'alter' a person's mental state, adding: 'It could promote psychosis or mood swings. 'It can alter your mental state and you might not know what you're doing.' The American-born Cambridge graduate (left) had been offered several US ambassadorial roles by Barack Obama in the weeks before her death and was planning a future with her new partner, David Shaw (right) Mrs Beerbower's estranged husband, whom she married in 1971, told the court that Mrs Beerbower had threatened suicide in the past. He also said she had previously 'over-medicated' and had mental health problems dating back to the 1970s. Mr Beebower told the inquest: 'I think there were some issues that go all the way back to the middle of the 70s. It was between two and three decades that she had been seeing people. ' He said that she became suicidal on occasions, adding: 'She would become very upset or angry and would make these kinds of statements and, for example, lock herself in the bathroom. 'As far as I'm aware she never undertook any action on any of these threats.' The inquest heard Mrs Beerbower was struggling to cope with her divorce from husband John and feared her family 'wished her dead'. Pictured: Mr Beerbower and the couple's daughter Sarah outside Westminster Coroner's Court . Despite her divorce not being finalised, Mrs Beerbower become engaged to David Shaw in May this year, the inquest heard. Mr Shaw described her as a 'highly intelligent women' and said they had started to make plans for the future. He said: 'We were planning for the next 20 years. We were planning about what countries we wanted to go, because she was being offered ambassadorial roles by President Obama.' Mrs Beerbower - one of the few female partners in a Wall Street law firm - had joined the first Clinton administration in 1993. She served as International Tax Counsel under Secretary Lloyd Bentsen before becoming Deputy Assistant Secretary for Tax Policy under Secretary Robert Rubin. She left US Government service in late 1996. Mrs Beerbower retired from the financial business in 2005 and moved to France with her daughter Sarah. Mrs Beerbower was found dead by her housekeeper at her apartment in Fulham, west London (pictured) She and her husband acquired an olive farm in Chateauneuf de Grasse, in the south of France in 2006. From 2012 onward, Cynthia divided her time between homes in Toronto, New York City, the South of France and London, the inquest heard. Home Office pathologist Dr Jon Van Der Walt gave the cause of death as Oxidon, Phetermine and paracetemol toxicity. He said the drugs could have been enough to create a loss of consciousness, or fatal arrhythmia. The coroner said she could not give a verdict of suicide because the drugs Mrs Beerbower was taking could have affected her decision-making. Dr Wilcox said: 'At the time of death she was under the influence of Phentermine - it may well have affected her ability to have clear intentions, or even think clearly in any sensible way.' She instead recorded a narrative verdict, adding: 'She died as a result of poly-drug overdose while under the influence of Phentermine.' For confidential support call the Samaritans in the UK on 08457 90 90 90, visit a local Samaritans branch or click here for details. परमाणु \\
\cline{1-9}
False & 15.0% & 57,944,719.3875 & \bfseries \color{red} 100.0000\% & 12.111 & 476.132 & 1,283 & 0.371 & High-flying tax lawter Cynthia Gibson Beerbower, 65, was found by her housekeeper at her apartment in Fulham, west London, in July this year . A high-flying tax lawyer who worked in the White House was found dead in the bath of her London home after taking an overdose while going through a 'painful' divorce. American-born Cynthia Gibson Beerbower, 65, was found by her housekeeper at her multi-million-pound apartment in Fulham, west London, on July 26 this year. The 'highly intelligent' former Wall Street lawyer had left three notes in her spare bedroom alongside a piece of card which read 'I can't do this anymore', an inquest has heard. The hearing was also told how the note included a request for Mrs Beerbower be buried in Cambridge, where she attended university. Mrs Beerbower had apparently been offered several ambassadorial roles by President Barack Obama in the weeks before her death and was planning a future with her new fiancé. But the hearing at Westminster Coroner's Court heard that Mrs Beerbower was struggling to cope with an ongoing divorce battle and feared her family 'wished her dead'. Her estranged husband also described how Mrs Beerbower had a history of mental illness and had threatened suicide several times. During the hearing, Coroner Fiona Wilcox told the court that a note left by the lawyer said: 'I can't do this any more, I hate you, I am sick'. The court also heard evidence from medical professionals who said Mrs Beerbower had started taking medication that she had not been prescribed. During the hearing, Dr Marios Pierides, a psychiatrist at the Nightingale Hospital in Marylebone, said Mrs Beerbower had visited him before her death. He described how she felt 'lonely' and 'isolated' while going through the process of divorcing her husband John Beerbower. He told the court: 'Cynthia was on various drugs at different time. 'When she came to me she was being treated with various medication, she was on anti-depressants and sleeping pills. 'She explained to me she was going through a very painful period of divorce. Also she was feeling quite lonely and isolated in her residence in Fulham. Dr Pierides added: 'She would say the family wished her dead. But she was, from my point of view, a life affirming person who enjoyed things until the awful circumstances around the divorce happened.' He continued: 'At the same time she was angry and lonely at her circumstances.' The court heard Mrs Beerbower had started taking Phentermine, an appetite suppressor, as well as Diazepam. Neither had been prescribed by the psychiatrist. Dr Pierides told the hearing that the drugs have the capability to 'alter' a person's mental state, adding: 'It could promote psychosis or mood swings. 'It can alter your mental state and you might not know what you're doing.' The American-born Cambridge graduate (left) had been offered several US ambassadorial roles by Barack Obama in the weeks before her death and was planning a future with her new partner, David Shaw (right) Mrs Beerbower's estranged husband, whom she married in 1971, told the court that Mrs Beerbower had threatened suicide in the past. He also said she had previously 'over-medicated' and had mental health problems dating back to the 1970s. Mr Beebower told the inquest: 'I think there were some issues that go all the way back to the middle of the 70s. It was between two and three decades that she had been seeing people. ' He said that she became suicidal on occasions, adding: 'She would become very upset or angry and would make these kinds of statements and, for example, lock herself in the bathroom. 'As far as I'm aware she never undertook any action on any of these threats.' The inquest heard Mrs Beerbower was struggling to cope with her divorce from husband John and feared her family 'wished her dead'. Pictured: Mr Beerbower and the couple's daughter Sarah outside Westminster Coroner's Court . Despite her divorce not being finalised, Mrs Beerbower become engaged to David Shaw in May this year, the inquest heard. Mr Shaw described her as a 'highly intelligent women' and said they had started to make plans for the future. He said: 'We were planning for the next 20 years. We were planning about what countries we wanted to go, because she was being offered ambassadorial roles by President Obama.' Mrs Beerbower - one of the few female partners in a Wall Street law firm - had joined the first Clinton administration in 1993. She served as International Tax Counsel under Secretary Lloyd Bentsen before becoming Deputy Assistant Secretary for Tax Policy under Secretary Robert Rubin. She left US Government service in late 1996. Mrs Beerbower retired from the financial business in 2005 and moved to France with her daughter Sarah. Mrs Beerbower was found dead by her housekeeper at her apartment in Fulham, west London (pictured) She and her husband acquired an olive farm in Chateauneuf de Grasse, in the south of France in 2006. From 2012 onward, Cynthia divided her time between homes in Toronto, New York City, the South of France and London, the inquest heard. Home Office pathologist Dr Jon Van Der Walt gave the cause of death as Oxidon, Phetermine and paracetemol toxicity. He said the drugs could have been enough to create a loss of consciousness, or fatal arrhythmia. The coroner said she could not give a verdict of suicide because the drugs Mrs Beerbower was taking could have affected her decision-making. Dr Wilcox said: 'At the time of death she was under the influence of Phentermine - it may well have affected her ability to have clear intentions, or even think clearly in any sensible way.' She instead recorded a narrative verdict, adding: 'She died as a result of poly-drug overdose while under the influence of Phentermine.' For confidential support call the Samaritans in the UK on 08457 90 90 90, visit a local Samaritans branch or click here for details.只要 \\
\cline{1-9}
False & 17.5% & 18,811,896.1195 & \bfseries \color{red} 100.0000\% & \bfseries \color{red} 11.946 & 477.936 & 1,283 & 0.373 & High-flying tax lawter Cynthia Gibson Beerbower, 65, was found by her housekeeper at her apartment in Fulham, west London, in July this year . A high-flying tax lawyer who worked in the White House was found dead in the bath of her London home after taking an overdose while going through a 'painful' divorce. American-born Cynthia Gibson Beerbower, 65, was found by her housekeeper at her multi-million-pound apartment in Fulham, west London, on July 26 this year. The 'highly intelligent' former Wall Street lawyer had left three notes in her spare bedroom alongside a piece of card which read 'I can't do this anymore', an inquest has heard. The hearing was also told how the note included a request for Mrs Beerbower be buried in Cambridge, where she attended university. Mrs Beerbower had apparently been offered several ambassadorial roles by President Barack Obama in the weeks before her death and was planning a future with her new fiancé. But the hearing at Westminster Coroner's Court heard that Mrs Beerbower was struggling to cope with an ongoing divorce battle and feared her family 'wished her dead'. Her estranged husband also described how Mrs Beerbower had a history of mental illness and had threatened suicide several times. During the hearing, Coroner Fiona Wilcox told the court that a note left by the lawyer said: 'I can't do this any more, I hate you, I am sick'. The court also heard evidence from medical professionals who said Mrs Beerbower had started taking medication that she had not been prescribed. During the hearing, Dr Marios Pierides, a psychiatrist at the Nightingale Hospital in Marylebone, said Mrs Beerbower had visited him before her death. He described how she felt 'lonely' and 'isolated' while going through the process of divorcing her husband John Beerbower. He told the court: 'Cynthia was on various drugs at different time. 'When she came to me she was being treated with various medication, she was on anti-depressants and sleeping pills. 'She explained to me she was going through a very painful period of divorce. Also she was feeling quite lonely and isolated in her residence in Fulham. Dr Pierides added: 'She would say the family wished her dead. But she was, from my point of view, a life affirming person who enjoyed things until the awful circumstances around the divorce happened.' He continued: 'At the same time she was angry and lonely at her circumstances.' The court heard Mrs Beerbower had started taking Phentermine, an appetite suppressor, as well as Diazepam. Neither had been prescribed by the psychiatrist. Dr Pierides told the hearing that the drugs have the capability to 'alter' a person's mental state, adding: 'It could promote psychosis or mood swings. 'It can alter your mental state and you might not know what you're doing.' The American-born Cambridge graduate (left) had been offered several US ambassadorial roles by Barack Obama in the weeks before her death and was planning a future with her new partner, David Shaw (right) Mrs Beerbower's estranged husband, whom she married in 1971, told the court that Mrs Beerbower had threatened suicide in the past. He also said she had previously 'over-medicated' and had mental health problems dating back to the 1970s. Mr Beebower told the inquest: 'I think there were some issues that go all the way back to the middle of the 70s. It was between two and three decades that she had been seeing people. ' He said that she became suicidal on occasions, adding: 'She would become very upset or angry and would make these kinds of statements and, for example, lock herself in the bathroom. 'As far as I'm aware she never undertook any action on any of these threats.' The inquest heard Mrs Beerbower was struggling to cope with her divorce from husband John and feared her family 'wished her dead'. Pictured: Mr Beerbower and the couple's daughter Sarah outside Westminster Coroner's Court . Despite her divorce not being finalised, Mrs Beerbower become engaged to David Shaw in May this year, the inquest heard. Mr Shaw described her as a 'highly intelligent women' and said they had started to make plans for the future. He said: 'We were planning for the next 20 years. We were planning about what countries we wanted to go, because she was being offered ambassadorial roles by President Obama.' Mrs Beerbower - one of the few female partners in a Wall Street law firm - had joined the first Clinton administration in 1993. She served as International Tax Counsel under Secretary Lloyd Bentsen before becoming Deputy Assistant Secretary for Tax Policy under Secretary Robert Rubin. She left US Government service in late 1996. Mrs Beerbower retired from the financial business in 2005 and moved to France with her daughter Sarah. Mrs Beerbower was found dead by her housekeeper at her apartment in Fulham, west London (pictured) She and her husband acquired an olive farm in Chateauneuf de Grasse, in the south of France in 2006. From 2012 onward, Cynthia divided her time between homes in Toronto, New York City, the South of France and London, the inquest heard. Home Office pathologist Dr Jon Van Der Walt gave the cause of death as Oxidon, Phetermine and paracetemol toxicity. He said the drugs could have been enough to create a loss of consciousness, or fatal arrhythmia. The coroner said she could not give a verdict of suicide because the drugs Mrs Beerbower was taking could have affected her decision-making. Dr Wilcox said: 'At the time of death she was under the influence of Phentermine - it may well have affected her ability to have clear intentions, or even think clearly in any sensible way.' She instead recorded a narrative verdict, adding: 'She died as a result of poly-drug overdose while under the influence of Phentermine.' For confidential support call the Samaritans in the UK on 08457 90 90 90, visit a local Samaritans branch or click here for details. Beverage \\
\cline{1-9}
False & 20.0% & 4,468,216.9751 & \bfseries \color{red} 100.0000\% & 12.123 & 476.763 & 1,283 & 0.372 & High-flying tax lawter Cynthia Gibson Beerbower, 65, was found by her housekeeper at her apartment in Fulham, west London, in July this year . A high-flying tax lawyer who worked in the White House was found dead in the bath of her London home after taking an overdose while going through a 'painful' divorce. American-born Cynthia Gibson Beerbower, 65, was found by her housekeeper at her multi-million-pound apartment in Fulham, west London, on July 26 this year. The 'highly intelligent' former Wall Street lawyer had left three notes in her spare bedroom alongside a piece of card which read 'I can't do this anymore', an inquest has heard. The hearing was also told how the note included a request for Mrs Beerbower be buried in Cambridge, where she attended university. Mrs Beerbower had apparently been offered several ambassadorial roles by President Barack Obama in the weeks before her death and was planning a future with her new fiancé. But the hearing at Westminster Coroner's Court heard that Mrs Beerbower was struggling to cope with an ongoing divorce battle and feared her family 'wished her dead'. Her estranged husband also described how Mrs Beerbower had a history of mental illness and had threatened suicide several times. During the hearing, Coroner Fiona Wilcox told the court that a note left by the lawyer said: 'I can't do this any more, I hate you, I am sick'. The court also heard evidence from medical professionals who said Mrs Beerbower had started taking medication that she had not been prescribed. During the hearing, Dr Marios Pierides, a psychiatrist at the Nightingale Hospital in Marylebone, said Mrs Beerbower had visited him before her death. He described how she felt 'lonely' and 'isolated' while going through the process of divorcing her husband John Beerbower. He told the court: 'Cynthia was on various drugs at different time. 'When she came to me she was being treated with various medication, she was on anti-depressants and sleeping pills. 'She explained to me she was going through a very painful period of divorce. Also she was feeling quite lonely and isolated in her residence in Fulham. Dr Pierides added: 'She would say the family wished her dead. But she was, from my point of view, a life affirming person who enjoyed things until the awful circumstances around the divorce happened.' He continued: 'At the same time she was angry and lonely at her circumstances.' The court heard Mrs Beerbower had started taking Phentermine, an appetite suppressor, as well as Diazepam. Neither had been prescribed by the psychiatrist. Dr Pierides told the hearing that the drugs have the capability to 'alter' a person's mental state, adding: 'It could promote psychosis or mood swings. 'It can alter your mental state and you might not know what you're doing.' The American-born Cambridge graduate (left) had been offered several US ambassadorial roles by Barack Obama in the weeks before her death and was planning a future with her new partner, David Shaw (right) Mrs Beerbower's estranged husband, whom she married in 1971, told the court that Mrs Beerbower had threatened suicide in the past. He also said she had previously 'over-medicated' and had mental health problems dating back to the 1970s. Mr Beebower told the inquest: 'I think there were some issues that go all the way back to the middle of the 70s. It was between two and three decades that she had been seeing people. ' He said that she became suicidal on occasions, adding: 'She would become very upset or angry and would make these kinds of statements and, for example, lock herself in the bathroom. 'As far as I'm aware she never undertook any action on any of these threats.' The inquest heard Mrs Beerbower was struggling to cope with her divorce from husband John and feared her family 'wished her dead'. Pictured: Mr Beerbower and the couple's daughter Sarah outside Westminster Coroner's Court . Despite her divorce not being finalised, Mrs Beerbower become engaged to David Shaw in May this year, the inquest heard. Mr Shaw described her as a 'highly intelligent women' and said they had started to make plans for the future. He said: 'We were planning for the next 20 years. We were planning about what countries we wanted to go, because she was being offered ambassadorial roles by President Obama.' Mrs Beerbower - one of the few female partners in a Wall Street law firm - had joined the first Clinton administration in 1993. She served as International Tax Counsel under Secretary Lloyd Bentsen before becoming Deputy Assistant Secretary for Tax Policy under Secretary Robert Rubin. She left US Government service in late 1996. Mrs Beerbower retired from the financial business in 2005 and moved to France with her daughter Sarah. Mrs Beerbower was found dead by her housekeeper at her apartment in Fulham, west London (pictured) She and her husband acquired an olive farm in Chateauneuf de Grasse, in the south of France in 2006. From 2012 onward, Cynthia divided her time between homes in Toronto, New York City, the South of France and London, the inquest heard. Home Office pathologist Dr Jon Van Der Walt gave the cause of death as Oxidon, Phetermine and paracetemol toxicity. He said the drugs could have been enough to create a loss of consciousness, or fatal arrhythmia. The coroner said she could not give a verdict of suicide because the drugs Mrs Beerbower was taking could have affected her decision-making. Dr Wilcox said: 'At the time of death she was under the influence of Phentermine - it may well have affected her ability to have clear intentions, or even think clearly in any sensible way.' She instead recorded a narrative verdict, adding: 'She died as a result of poly-drug overdose while under the influence of Phentermine.' For confidential support call the Samaritans in the UK on 08457 90 90 90, visit a local Samaritans branch or click here for details.ownika \\
\cline{1-9}
False & 22.5% & 12,929,214.5167 & \bfseries \color{red} 100.0000\% & 13.409 & 476.553 & 1,283 & 0.371 & High-flying tax lawter Cynthia Gibson Beerbower, 65, was found by her housekeeper at her apartment in Fulham, west London, in July this year . A high-flying tax lawyer who worked in the White House was found dead in the bath of her London home after taking an overdose while going through a 'painful' divorce. American-born Cynthia Gibson Beerbower, 65, was found by her housekeeper at her multi-million-pound apartment in Fulham, west London, on July 26 this year. The 'highly intelligent' former Wall Street lawyer had left three notes in her spare bedroom alongside a piece of card which read 'I can't do this anymore', an inquest has heard. The hearing was also told how the note included a request for Mrs Beerbower be buried in Cambridge, where she attended university. Mrs Beerbower had apparently been offered several ambassadorial roles by President Barack Obama in the weeks before her death and was planning a future with her new fiancé. But the hearing at Westminster Coroner's Court heard that Mrs Beerbower was struggling to cope with an ongoing divorce battle and feared her family 'wished her dead'. Her estranged husband also described how Mrs Beerbower had a history of mental illness and had threatened suicide several times. During the hearing, Coroner Fiona Wilcox told the court that a note left by the lawyer said: 'I can't do this any more, I hate you, I am sick'. The court also heard evidence from medical professionals who said Mrs Beerbower had started taking medication that she had not been prescribed. During the hearing, Dr Marios Pierides, a psychiatrist at the Nightingale Hospital in Marylebone, said Mrs Beerbower had visited him before her death. He described how she felt 'lonely' and 'isolated' while going through the process of divorcing her husband John Beerbower. He told the court: 'Cynthia was on various drugs at different time. 'When she came to me she was being treated with various medication, she was on anti-depressants and sleeping pills. 'She explained to me she was going through a very painful period of divorce. Also she was feeling quite lonely and isolated in her residence in Fulham. Dr Pierides added: 'She would say the family wished her dead. But she was, from my point of view, a life affirming person who enjoyed things until the awful circumstances around the divorce happened.' He continued: 'At the same time she was angry and lonely at her circumstances.' The court heard Mrs Beerbower had started taking Phentermine, an appetite suppressor, as well as Diazepam. Neither had been prescribed by the psychiatrist. Dr Pierides told the hearing that the drugs have the capability to 'alter' a person's mental state, adding: 'It could promote psychosis or mood swings. 'It can alter your mental state and you might not know what you're doing.' The American-born Cambridge graduate (left) had been offered several US ambassadorial roles by Barack Obama in the weeks before her death and was planning a future with her new partner, David Shaw (right) Mrs Beerbower's estranged husband, whom she married in 1971, told the court that Mrs Beerbower had threatened suicide in the past. He also said she had previously 'over-medicated' and had mental health problems dating back to the 1970s. Mr Beebower told the inquest: 'I think there were some issues that go all the way back to the middle of the 70s. It was between two and three decades that she had been seeing people. ' He said that she became suicidal on occasions, adding: 'She would become very upset or angry and would make these kinds of statements and, for example, lock herself in the bathroom. 'As far as I'm aware she never undertook any action on any of these threats.' The inquest heard Mrs Beerbower was struggling to cope with her divorce from husband John and feared her family 'wished her dead'. Pictured: Mr Beerbower and the couple's daughter Sarah outside Westminster Coroner's Court . Despite her divorce not being finalised, Mrs Beerbower become engaged to David Shaw in May this year, the inquest heard. Mr Shaw described her as a 'highly intelligent women' and said they had started to make plans for the future. He said: 'We were planning for the next 20 years. We were planning about what countries we wanted to go, because she was being offered ambassadorial roles by President Obama.' Mrs Beerbower - one of the few female partners in a Wall Street law firm - had joined the first Clinton administration in 1993. She served as International Tax Counsel under Secretary Lloyd Bentsen before becoming Deputy Assistant Secretary for Tax Policy under Secretary Robert Rubin. She left US Government service in late 1996. Mrs Beerbower retired from the financial business in 2005 and moved to France with her daughter Sarah. Mrs Beerbower was found dead by her housekeeper at her apartment in Fulham, west London (pictured) She and her husband acquired an olive farm in Chateauneuf de Grasse, in the south of France in 2006. From 2012 onward, Cynthia divided her time between homes in Toronto, New York City, the South of France and London, the inquest heard. Home Office pathologist Dr Jon Van Der Walt gave the cause of death as Oxidon, Phetermine and paracetemol toxicity. He said the drugs could have been enough to create a loss of consciousness, or fatal arrhythmia. The coroner said she could not give a verdict of suicide because the drugs Mrs Beerbower was taking could have affected her decision-making. Dr Wilcox said: 'At the time of death she was under the influence of Phentermine - it may well have affected her ability to have clear intentions, or even think clearly in any sensible way.' She instead recorded a narrative verdict, adding: 'She died as a result of poly-drug overdose while under the influence of Phentermine.' For confidential support call the Samaritans in the UK on 08457 90 90 90, visit a local Samaritans branch or click here for details.ettle \\
\cline{1-9}
False & 25.0% & 11,409,991.7638 & \bfseries \color{red} 100.0000\% & 12.018 & 474.906 & 1,283 & 0.370 & High-flying tax lawter Cynthia Gibson Beerbower, 65, was found by her housekeeper at her apartment in Fulham, west London, in July this year . A high-flying tax lawyer who worked in the White House was found dead in the bath of her London home after taking an overdose while going through a 'painful' divorce. American-born Cynthia Gibson Beerbower, 65, was found by her housekeeper at her multi-million-pound apartment in Fulham, west London, on July 26 this year. The 'highly intelligent' former Wall Street lawyer had left three notes in her spare bedroom alongside a piece of card which read 'I can't do this anymore', an inquest has heard. The hearing was also told how the note included a request for Mrs Beerbower be buried in Cambridge, where she attended university. Mrs Beerbower had apparently been offered several ambassadorial roles by President Barack Obama in the weeks before her death and was planning a future with her new fiancé. But the hearing at Westminster Coroner's Court heard that Mrs Beerbower was struggling to cope with an ongoing divorce battle and feared her family 'wished her dead'. Her estranged husband also described how Mrs Beerbower had a history of mental illness and had threatened suicide several times. During the hearing, Coroner Fiona Wilcox told the court that a note left by the lawyer said: 'I can't do this any more, I hate you, I am sick'. The court also heard evidence from medical professionals who said Mrs Beerbower had started taking medication that she had not been prescribed. During the hearing, Dr Marios Pierides, a psychiatrist at the Nightingale Hospital in Marylebone, said Mrs Beerbower had visited him before her death. He described how she felt 'lonely' and 'isolated' while going through the process of divorcing her husband John Beerbower. He told the court: 'Cynthia was on various drugs at different time. 'When she came to me she was being treated with various medication, she was on anti-depressants and sleeping pills. 'She explained to me she was going through a very painful period of divorce. Also she was feeling quite lonely and isolated in her residence in Fulham. Dr Pierides added: 'She would say the family wished her dead. But she was, from my point of view, a life affirming person who enjoyed things until the awful circumstances around the divorce happened.' He continued: 'At the same time she was angry and lonely at her circumstances.' The court heard Mrs Beerbower had started taking Phentermine, an appetite suppressor, as well as Diazepam. Neither had been prescribed by the psychiatrist. Dr Pierides told the hearing that the drugs have the capability to 'alter' a person's mental state, adding: 'It could promote psychosis or mood swings. 'It can alter your mental state and you might not know what you're doing.' The American-born Cambridge graduate (left) had been offered several US ambassadorial roles by Barack Obama in the weeks before her death and was planning a future with her new partner, David Shaw (right) Mrs Beerbower's estranged husband, whom she married in 1971, told the court that Mrs Beerbower had threatened suicide in the past. He also said she had previously 'over-medicated' and had mental health problems dating back to the 1970s. Mr Beebower told the inquest: 'I think there were some issues that go all the way back to the middle of the 70s. It was between two and three decades that she had been seeing people. ' He said that she became suicidal on occasions, adding: 'She would become very upset or angry and would make these kinds of statements and, for example, lock herself in the bathroom. 'As far as I'm aware she never undertook any action on any of these threats.' The inquest heard Mrs Beerbower was struggling to cope with her divorce from husband John and feared her family 'wished her dead'. Pictured: Mr Beerbower and the couple's daughter Sarah outside Westminster Coroner's Court . Despite her divorce not being finalised, Mrs Beerbower become engaged to David Shaw in May this year, the inquest heard. Mr Shaw described her as a 'highly intelligent women' and said they had started to make plans for the future. He said: 'We were planning for the next 20 years. We were planning about what countries we wanted to go, because she was being offered ambassadorial roles by President Obama.' Mrs Beerbower - one of the few female partners in a Wall Street law firm - had joined the first Clinton administration in 1993. She served as International Tax Counsel under Secretary Lloyd Bentsen before becoming Deputy Assistant Secretary for Tax Policy under Secretary Robert Rubin. She left US Government service in late 1996. Mrs Beerbower retired from the financial business in 2005 and moved to France with her daughter Sarah. Mrs Beerbower was found dead by her housekeeper at her apartment in Fulham, west London (pictured) She and her husband acquired an olive farm in Chateauneuf de Grasse, in the south of France in 2006. From 2012 onward, Cynthia divided her time between homes in Toronto, New York City, the South of France and London, the inquest heard. Home Office pathologist Dr Jon Van Der Walt gave the cause of death as Oxidon, Phetermine and paracetemol toxicity. He said the drugs could have been enough to create a loss of consciousness, or fatal arrhythmia. The coroner said she could not give a verdict of suicide because the drugs Mrs Beerbower was taking could have affected her decision-making. Dr Wilcox said: 'At the time of death she was under the influence of Phentermine - it may well have affected her ability to have clear intentions, or even think clearly in any sensible way.' She instead recorded a narrative verdict, adding: 'She died as a result of poly-drug overdose while under the influence of Phentermine.' For confidential support call the Samaritans in the UK on 08457 90 90 90, visit a local Samaritans branch or click here for details.ঞ্জে \\
\cline{1-9}
False & 27.5% & 14,650,719.4290 & \bfseries \color{red} 100.0000\% & 13.172 & 476.408 & 1,283 & 0.371 & High-flying tax lawter Cynthia Gibson Beerbower, 65, was found by her housekeeper at her apartment in Fulham, west London, in July this year . A high-flying tax lawyer who worked in the White House was found dead in the bath of her London home after taking an overdose while going through a 'painful' divorce. American-born Cynthia Gibson Beerbower, 65, was found by her housekeeper at her multi-million-pound apartment in Fulham, west London, on July 26 this year. The 'highly intelligent' former Wall Street lawyer had left three notes in her spare bedroom alongside a piece of card which read 'I can't do this anymore', an inquest has heard. The hearing was also told how the note included a request for Mrs Beerbower be buried in Cambridge, where she attended university. Mrs Beerbower had apparently been offered several ambassadorial roles by President Barack Obama in the weeks before her death and was planning a future with her new fiancé. But the hearing at Westminster Coroner's Court heard that Mrs Beerbower was struggling to cope with an ongoing divorce battle and feared her family 'wished her dead'. Her estranged husband also described how Mrs Beerbower had a history of mental illness and had threatened suicide several times. During the hearing, Coroner Fiona Wilcox told the court that a note left by the lawyer said: 'I can't do this any more, I hate you, I am sick'. The court also heard evidence from medical professionals who said Mrs Beerbower had started taking medication that she had not been prescribed. During the hearing, Dr Marios Pierides, a psychiatrist at the Nightingale Hospital in Marylebone, said Mrs Beerbower had visited him before her death. He described how she felt 'lonely' and 'isolated' while going through the process of divorcing her husband John Beerbower. He told the court: 'Cynthia was on various drugs at different time. 'When she came to me she was being treated with various medication, she was on anti-depressants and sleeping pills. 'She explained to me she was going through a very painful period of divorce. Also she was feeling quite lonely and isolated in her residence in Fulham. Dr Pierides added: 'She would say the family wished her dead. But she was, from my point of view, a life affirming person who enjoyed things until the awful circumstances around the divorce happened.' He continued: 'At the same time she was angry and lonely at her circumstances.' The court heard Mrs Beerbower had started taking Phentermine, an appetite suppressor, as well as Diazepam. Neither had been prescribed by the psychiatrist. Dr Pierides told the hearing that the drugs have the capability to 'alter' a person's mental state, adding: 'It could promote psychosis or mood swings. 'It can alter your mental state and you might not know what you're doing.' The American-born Cambridge graduate (left) had been offered several US ambassadorial roles by Barack Obama in the weeks before her death and was planning a future with her new partner, David Shaw (right) Mrs Beerbower's estranged husband, whom she married in 1971, told the court that Mrs Beerbower had threatened suicide in the past. He also said she had previously 'over-medicated' and had mental health problems dating back to the 1970s. Mr Beebower told the inquest: 'I think there were some issues that go all the way back to the middle of the 70s. It was between two and three decades that she had been seeing people. ' He said that she became suicidal on occasions, adding: 'She would become very upset or angry and would make these kinds of statements and, for example, lock herself in the bathroom. 'As far as I'm aware she never undertook any action on any of these threats.' The inquest heard Mrs Beerbower was struggling to cope with her divorce from husband John and feared her family 'wished her dead'. Pictured: Mr Beerbower and the couple's daughter Sarah outside Westminster Coroner's Court . Despite her divorce not being finalised, Mrs Beerbower become engaged to David Shaw in May this year, the inquest heard. Mr Shaw described her as a 'highly intelligent women' and said they had started to make plans for the future. He said: 'We were planning for the next 20 years. We were planning about what countries we wanted to go, because she was being offered ambassadorial roles by President Obama.' Mrs Beerbower - one of the few female partners in a Wall Street law firm - had joined the first Clinton administration in 1993. She served as International Tax Counsel under Secretary Lloyd Bentsen before becoming Deputy Assistant Secretary for Tax Policy under Secretary Robert Rubin. She left US Government service in late 1996. Mrs Beerbower retired from the financial business in 2005 and moved to France with her daughter Sarah. Mrs Beerbower was found dead by her housekeeper at her apartment in Fulham, west London (pictured) She and her husband acquired an olive farm in Chateauneuf de Grasse, in the south of France in 2006. From 2012 onward, Cynthia divided her time between homes in Toronto, New York City, the South of France and London, the inquest heard. Home Office pathologist Dr Jon Van Der Walt gave the cause of death as Oxidon, Phetermine and paracetemol toxicity. He said the drugs could have been enough to create a loss of consciousness, or fatal arrhythmia. The coroner said she could not give a verdict of suicide because the drugs Mrs Beerbower was taking could have affected her decision-making. Dr Wilcox said: 'At the time of death she was under the influence of Phentermine - it may well have affected her ability to have clear intentions, or even think clearly in any sensible way.' She instead recorded a narrative verdict, adding: 'She died as a result of poly-drug overdose while under the influence of Phentermine.' For confidential support call the Samaritans in the UK on 08457 90 90 90, visit a local Samaritans branch or click here for details. Sparta \\
\cline{1-9}
False & 30.0% & 12,929,214.5167 & \bfseries \color{red} 100.0000\% & 12.016 & 481.040 & 1,283 & 0.375 & High-flying tax lawter Cynthia Gibson Beerbower, 65, was found by her housekeeper at her apartment in Fulham, west London, in July this year . A high-flying tax lawyer who worked in the White House was found dead in the bath of her London home after taking an overdose while going through a 'painful' divorce. American-born Cynthia Gibson Beerbower, 65, was found by her housekeeper at her multi-million-pound apartment in Fulham, west London, on July 26 this year. The 'highly intelligent' former Wall Street lawyer had left three notes in her spare bedroom alongside a piece of card which read 'I can't do this anymore', an inquest has heard. The hearing was also told how the note included a request for Mrs Beerbower be buried in Cambridge, where she attended university. Mrs Beerbower had apparently been offered several ambassadorial roles by President Barack Obama in the weeks before her death and was planning a future with her new fiancé. But the hearing at Westminster Coroner's Court heard that Mrs Beerbower was struggling to cope with an ongoing divorce battle and feared her family 'wished her dead'. Her estranged husband also described how Mrs Beerbower had a history of mental illness and had threatened suicide several times. During the hearing, Coroner Fiona Wilcox told the court that a note left by the lawyer said: 'I can't do this any more, I hate you, I am sick'. The court also heard evidence from medical professionals who said Mrs Beerbower had started taking medication that she had not been prescribed. During the hearing, Dr Marios Pierides, a psychiatrist at the Nightingale Hospital in Marylebone, said Mrs Beerbower had visited him before her death. He described how she felt 'lonely' and 'isolated' while going through the process of divorcing her husband John Beerbower. He told the court: 'Cynthia was on various drugs at different time. 'When she came to me she was being treated with various medication, she was on anti-depressants and sleeping pills. 'She explained to me she was going through a very painful period of divorce. Also she was feeling quite lonely and isolated in her residence in Fulham. Dr Pierides added: 'She would say the family wished her dead. But she was, from my point of view, a life affirming person who enjoyed things until the awful circumstances around the divorce happened.' He continued: 'At the same time she was angry and lonely at her circumstances.' The court heard Mrs Beerbower had started taking Phentermine, an appetite suppressor, as well as Diazepam. Neither had been prescribed by the psychiatrist. Dr Pierides told the hearing that the drugs have the capability to 'alter' a person's mental state, adding: 'It could promote psychosis or mood swings. 'It can alter your mental state and you might not know what you're doing.' The American-born Cambridge graduate (left) had been offered several US ambassadorial roles by Barack Obama in the weeks before her death and was planning a future with her new partner, David Shaw (right) Mrs Beerbower's estranged husband, whom she married in 1971, told the court that Mrs Beerbower had threatened suicide in the past. He also said she had previously 'over-medicated' and had mental health problems dating back to the 1970s. Mr Beebower told the inquest: 'I think there were some issues that go all the way back to the middle of the 70s. It was between two and three decades that she had been seeing people. ' He said that she became suicidal on occasions, adding: 'She would become very upset or angry and would make these kinds of statements and, for example, lock herself in the bathroom. 'As far as I'm aware she never undertook any action on any of these threats.' The inquest heard Mrs Beerbower was struggling to cope with her divorce from husband John and feared her family 'wished her dead'. Pictured: Mr Beerbower and the couple's daughter Sarah outside Westminster Coroner's Court . Despite her divorce not being finalised, Mrs Beerbower become engaged to David Shaw in May this year, the inquest heard. Mr Shaw described her as a 'highly intelligent women' and said they had started to make plans for the future. He said: 'We were planning for the next 20 years. We were planning about what countries we wanted to go, because she was being offered ambassadorial roles by President Obama.' Mrs Beerbower - one of the few female partners in a Wall Street law firm - had joined the first Clinton administration in 1993. She served as International Tax Counsel under Secretary Lloyd Bentsen before becoming Deputy Assistant Secretary for Tax Policy under Secretary Robert Rubin. She left US Government service in late 1996. Mrs Beerbower retired from the financial business in 2005 and moved to France with her daughter Sarah. Mrs Beerbower was found dead by her housekeeper at her apartment in Fulham, west London (pictured) She and her husband acquired an olive farm in Chateauneuf de Grasse, in the south of France in 2006. From 2012 onward, Cynthia divided her time between homes in Toronto, New York City, the South of France and London, the inquest heard. Home Office pathologist Dr Jon Van Der Walt gave the cause of death as Oxidon, Phetermine and paracetemol toxicity. He said the drugs could have been enough to create a loss of consciousness, or fatal arrhythmia. The coroner said she could not give a verdict of suicide because the drugs Mrs Beerbower was taking could have affected her decision-making. Dr Wilcox said: 'At the time of death she was under the influence of Phentermine - it may well have affected her ability to have clear intentions, or even think clearly in any sensible way.' She instead recorded a narrative verdict, adding: 'She died as a result of poly-drug overdose while under the influence of Phentermine.' For confidential support call the Samaritans in the UK on 08457 90 90 90, visit a local Samaritans branch or click here for details. prostate \\
\cline{1-9}
False & 32.5% & 16,601,440.0572 & \bfseries \color{red} 100.0000\% & 12.072 & 477.377 & 1,283 & 0.372 & High-flying tax lawter Cynthia Gibson Beerbower, 65, was found by her housekeeper at her apartment in Fulham, west London, in July this year . A high-flying tax lawyer who worked in the White House was found dead in the bath of her London home after taking an overdose while going through a 'painful' divorce. American-born Cynthia Gibson Beerbower, 65, was found by her housekeeper at her multi-million-pound apartment in Fulham, west London, on July 26 this year. The 'highly intelligent' former Wall Street lawyer had left three notes in her spare bedroom alongside a piece of card which read 'I can't do this anymore', an inquest has heard. The hearing was also told how the note included a request for Mrs Beerbower be buried in Cambridge, where she attended university. Mrs Beerbower had apparently been offered several ambassadorial roles by President Barack Obama in the weeks before her death and was planning a future with her new fiancé. But the hearing at Westminster Coroner's Court heard that Mrs Beerbower was struggling to cope with an ongoing divorce battle and feared her family 'wished her dead'. Her estranged husband also described how Mrs Beerbower had a history of mental illness and had threatened suicide several times. During the hearing, Coroner Fiona Wilcox told the court that a note left by the lawyer said: 'I can't do this any more, I hate you, I am sick'. The court also heard evidence from medical professionals who said Mrs Beerbower had started taking medication that she had not been prescribed. During the hearing, Dr Marios Pierides, a psychiatrist at the Nightingale Hospital in Marylebone, said Mrs Beerbower had visited him before her death. He described how she felt 'lonely' and 'isolated' while going through the process of divorcing her husband John Beerbower. He told the court: 'Cynthia was on various drugs at different time. 'When she came to me she was being treated with various medication, she was on anti-depressants and sleeping pills. 'She explained to me she was going through a very painful period of divorce. Also she was feeling quite lonely and isolated in her residence in Fulham. Dr Pierides added: 'She would say the family wished her dead. But she was, from my point of view, a life affirming person who enjoyed things until the awful circumstances around the divorce happened.' He continued: 'At the same time she was angry and lonely at her circumstances.' The court heard Mrs Beerbower had started taking Phentermine, an appetite suppressor, as well as Diazepam. Neither had been prescribed by the psychiatrist. Dr Pierides told the hearing that the drugs have the capability to 'alter' a person's mental state, adding: 'It could promote psychosis or mood swings. 'It can alter your mental state and you might not know what you're doing.' The American-born Cambridge graduate (left) had been offered several US ambassadorial roles by Barack Obama in the weeks before her death and was planning a future with her new partner, David Shaw (right) Mrs Beerbower's estranged husband, whom she married in 1971, told the court that Mrs Beerbower had threatened suicide in the past. He also said she had previously 'over-medicated' and had mental health problems dating back to the 1970s. Mr Beebower told the inquest: 'I think there were some issues that go all the way back to the middle of the 70s. It was between two and three decades that she had been seeing people. ' He said that she became suicidal on occasions, adding: 'She would become very upset or angry and would make these kinds of statements and, for example, lock herself in the bathroom. 'As far as I'm aware she never undertook any action on any of these threats.' The inquest heard Mrs Beerbower was struggling to cope with her divorce from husband John and feared her family 'wished her dead'. Pictured: Mr Beerbower and the couple's daughter Sarah outside Westminster Coroner's Court . Despite her divorce not being finalised, Mrs Beerbower become engaged to David Shaw in May this year, the inquest heard. Mr Shaw described her as a 'highly intelligent women' and said they had started to make plans for the future. He said: 'We were planning for the next 20 years. We were planning about what countries we wanted to go, because she was being offered ambassadorial roles by President Obama.' Mrs Beerbower - one of the few female partners in a Wall Street law firm - had joined the first Clinton administration in 1993. She served as International Tax Counsel under Secretary Lloyd Bentsen before becoming Deputy Assistant Secretary for Tax Policy under Secretary Robert Rubin. She left US Government service in late 1996. Mrs Beerbower retired from the financial business in 2005 and moved to France with her daughter Sarah. Mrs Beerbower was found dead by her housekeeper at her apartment in Fulham, west London (pictured) She and her husband acquired an olive farm in Chateauneuf de Grasse, in the south of France in 2006. From 2012 onward, Cynthia divided her time between homes in Toronto, New York City, the South of France and London, the inquest heard. Home Office pathologist Dr Jon Van Der Walt gave the cause of death as Oxidon, Phetermine and paracetemol toxicity. He said the drugs could have been enough to create a loss of consciousness, or fatal arrhythmia. The coroner said she could not give a verdict of suicide because the drugs Mrs Beerbower was taking could have affected her decision-making. Dr Wilcox said: 'At the time of death she was under the influence of Phentermine - it may well have affected her ability to have clear intentions, or even think clearly in any sensible way.' She instead recorded a narrative verdict, adding: 'She died as a result of poly-drug overdose while under the influence of Phentermine.' For confidential support call the Samaritans in the UK on 08457 90 90 90, visit a local Samaritans branch or click here for details. pât \\
\cline{1-9}
False & 35.0% & 27,371,147.3466 & \bfseries \color{red} 100.0000\% & 12.178 & 478.582 & 1,283 & 0.373 & High-flying tax lawter Cynthia Gibson Beerbower, 65, was found by her housekeeper at her apartment in Fulham, west London, in July this year . A high-flying tax lawyer who worked in the White House was found dead in the bath of her London home after taking an overdose while going through a 'painful' divorce. American-born Cynthia Gibson Beerbower, 65, was found by her housekeeper at her multi-million-pound apartment in Fulham, west London, on July 26 this year. The 'highly intelligent' former Wall Street lawyer had left three notes in her spare bedroom alongside a piece of card which read 'I can't do this anymore', an inquest has heard. The hearing was also told how the note included a request for Mrs Beerbower be buried in Cambridge, where she attended university. Mrs Beerbower had apparently been offered several ambassadorial roles by President Barack Obama in the weeks before her death and was planning a future with her new fiancé. But the hearing at Westminster Coroner's Court heard that Mrs Beerbower was struggling to cope with an ongoing divorce battle and feared her family 'wished her dead'. Her estranged husband also described how Mrs Beerbower had a history of mental illness and had threatened suicide several times. During the hearing, Coroner Fiona Wilcox told the court that a note left by the lawyer said: 'I can't do this any more, I hate you, I am sick'. The court also heard evidence from medical professionals who said Mrs Beerbower had started taking medication that she had not been prescribed. During the hearing, Dr Marios Pierides, a psychiatrist at the Nightingale Hospital in Marylebone, said Mrs Beerbower had visited him before her death. He described how she felt 'lonely' and 'isolated' while going through the process of divorcing her husband John Beerbower. He told the court: 'Cynthia was on various drugs at different time. 'When she came to me she was being treated with various medication, she was on anti-depressants and sleeping pills. 'She explained to me she was going through a very painful period of divorce. Also she was feeling quite lonely and isolated in her residence in Fulham. Dr Pierides added: 'She would say the family wished her dead. But she was, from my point of view, a life affirming person who enjoyed things until the awful circumstances around the divorce happened.' He continued: 'At the same time she was angry and lonely at her circumstances.' The court heard Mrs Beerbower had started taking Phentermine, an appetite suppressor, as well as Diazepam. Neither had been prescribed by the psychiatrist. Dr Pierides told the hearing that the drugs have the capability to 'alter' a person's mental state, adding: 'It could promote psychosis or mood swings. 'It can alter your mental state and you might not know what you're doing.' The American-born Cambridge graduate (left) had been offered several US ambassadorial roles by Barack Obama in the weeks before her death and was planning a future with her new partner, David Shaw (right) Mrs Beerbower's estranged husband, whom she married in 1971, told the court that Mrs Beerbower had threatened suicide in the past. He also said she had previously 'over-medicated' and had mental health problems dating back to the 1970s. Mr Beebower told the inquest: 'I think there were some issues that go all the way back to the middle of the 70s. It was between two and three decades that she had been seeing people. ' He said that she became suicidal on occasions, adding: 'She would become very upset or angry and would make these kinds of statements and, for example, lock herself in the bathroom. 'As far as I'm aware she never undertook any action on any of these threats.' The inquest heard Mrs Beerbower was struggling to cope with her divorce from husband John and feared her family 'wished her dead'. Pictured: Mr Beerbower and the couple's daughter Sarah outside Westminster Coroner's Court . Despite her divorce not being finalised, Mrs Beerbower become engaged to David Shaw in May this year, the inquest heard. Mr Shaw described her as a 'highly intelligent women' and said they had started to make plans for the future. He said: 'We were planning for the next 20 years. We were planning about what countries we wanted to go, because she was being offered ambassadorial roles by President Obama.' Mrs Beerbower - one of the few female partners in a Wall Street law firm - had joined the first Clinton administration in 1993. She served as International Tax Counsel under Secretary Lloyd Bentsen before becoming Deputy Assistant Secretary for Tax Policy under Secretary Robert Rubin. She left US Government service in late 1996. Mrs Beerbower retired from the financial business in 2005 and moved to France with her daughter Sarah. Mrs Beerbower was found dead by her housekeeper at her apartment in Fulham, west London (pictured) She and her husband acquired an olive farm in Chateauneuf de Grasse, in the south of France in 2006. From 2012 onward, Cynthia divided her time between homes in Toronto, New York City, the South of France and London, the inquest heard. Home Office pathologist Dr Jon Van Der Walt gave the cause of death as Oxidon, Phetermine and paracetemol toxicity. He said the drugs could have been enough to create a loss of consciousness, or fatal arrhythmia. The coroner said she could not give a verdict of suicide because the drugs Mrs Beerbower was taking could have affected her decision-making. Dr Wilcox said: 'At the time of death she was under the influence of Phentermine - it may well have affected her ability to have clear intentions, or even think clearly in any sensible way.' She instead recorded a narrative verdict, adding: 'She died as a result of poly-drug overdose while under the influence of Phentermine.' For confidential support call the Samaritans in the UK on 08457 90 90 90, visit a local Samaritans branch or click here for details.wiąz \\
\cline{1-9}
False & 37.5% & 16,601,440.0572 & \bfseries \color{red} 100.0000\% & 12.023 & 480.957 & 1,283 & 0.375 & High-flying tax lawter Cynthia Gibson Beerbower, 65, was found by her housekeeper at her apartment in Fulham, west London, in July this year . A high-flying tax lawyer who worked in the White House was found dead in the bath of her London home after taking an overdose while going through a 'painful' divorce. American-born Cynthia Gibson Beerbower, 65, was found by her housekeeper at her multi-million-pound apartment in Fulham, west London, on July 26 this year. The 'highly intelligent' former Wall Street lawyer had left three notes in her spare bedroom alongside a piece of card which read 'I can't do this anymore', an inquest has heard. The hearing was also told how the note included a request for Mrs Beerbower be buried in Cambridge, where she attended university. Mrs Beerbower had apparently been offered several ambassadorial roles by President Barack Obama in the weeks before her death and was planning a future with her new fiancé. But the hearing at Westminster Coroner's Court heard that Mrs Beerbower was struggling to cope with an ongoing divorce battle and feared her family 'wished her dead'. Her estranged husband also described how Mrs Beerbower had a history of mental illness and had threatened suicide several times. During the hearing, Coroner Fiona Wilcox told the court that a note left by the lawyer said: 'I can't do this any more, I hate you, I am sick'. The court also heard evidence from medical professionals who said Mrs Beerbower had started taking medication that she had not been prescribed. During the hearing, Dr Marios Pierides, a psychiatrist at the Nightingale Hospital in Marylebone, said Mrs Beerbower had visited him before her death. He described how she felt 'lonely' and 'isolated' while going through the process of divorcing her husband John Beerbower. He told the court: 'Cynthia was on various drugs at different time. 'When she came to me she was being treated with various medication, she was on anti-depressants and sleeping pills. 'She explained to me she was going through a very painful period of divorce. Also she was feeling quite lonely and isolated in her residence in Fulham. Dr Pierides added: 'She would say the family wished her dead. But she was, from my point of view, a life affirming person who enjoyed things until the awful circumstances around the divorce happened.' He continued: 'At the same time she was angry and lonely at her circumstances.' The court heard Mrs Beerbower had started taking Phentermine, an appetite suppressor, as well as Diazepam. Neither had been prescribed by the psychiatrist. Dr Pierides told the hearing that the drugs have the capability to 'alter' a person's mental state, adding: 'It could promote psychosis or mood swings. 'It can alter your mental state and you might not know what you're doing.' The American-born Cambridge graduate (left) had been offered several US ambassadorial roles by Barack Obama in the weeks before her death and was planning a future with her new partner, David Shaw (right) Mrs Beerbower's estranged husband, whom she married in 1971, told the court that Mrs Beerbower had threatened suicide in the past. He also said she had previously 'over-medicated' and had mental health problems dating back to the 1970s. Mr Beebower told the inquest: 'I think there were some issues that go all the way back to the middle of the 70s. It was between two and three decades that she had been seeing people. ' He said that she became suicidal on occasions, adding: 'She would become very upset or angry and would make these kinds of statements and, for example, lock herself in the bathroom. 'As far as I'm aware she never undertook any action on any of these threats.' The inquest heard Mrs Beerbower was struggling to cope with her divorce from husband John and feared her family 'wished her dead'. Pictured: Mr Beerbower and the couple's daughter Sarah outside Westminster Coroner's Court . Despite her divorce not being finalised, Mrs Beerbower become engaged to David Shaw in May this year, the inquest heard. Mr Shaw described her as a 'highly intelligent women' and said they had started to make plans for the future. He said: 'We were planning for the next 20 years. We were planning about what countries we wanted to go, because she was being offered ambassadorial roles by President Obama.' Mrs Beerbower - one of the few female partners in a Wall Street law firm - had joined the first Clinton administration in 1993. She served as International Tax Counsel under Secretary Lloyd Bentsen before becoming Deputy Assistant Secretary for Tax Policy under Secretary Robert Rubin. She left US Government service in late 1996. Mrs Beerbower retired from the financial business in 2005 and moved to France with her daughter Sarah. Mrs Beerbower was found dead by her housekeeper at her apartment in Fulham, west London (pictured) She and her husband acquired an olive farm in Chateauneuf de Grasse, in the south of France in 2006. From 2012 onward, Cynthia divided her time between homes in Toronto, New York City, the South of France and London, the inquest heard. Home Office pathologist Dr Jon Van Der Walt gave the cause of death as Oxidon, Phetermine and paracetemol toxicity. He said the drugs could have been enough to create a loss of consciousness, or fatal arrhythmia. The coroner said she could not give a verdict of suicide because the drugs Mrs Beerbower was taking could have affected her decision-making. Dr Wilcox said: 'At the time of death she was under the influence of Phentermine - it may well have affected her ability to have clear intentions, or even think clearly in any sensible way.' She instead recorded a narrative verdict, adding: 'She died as a result of poly-drug overdose while under the influence of Phentermine.' For confidential support call the Samaritans in the UK on 08457 90 90 90, visit a local Samaritans branch or click here for details. austerity \\
\cline{1-9}
False & 40.0% & 27,371,147.3466 & \bfseries \color{red} 100.0000\% & 12.054 & 476.401 & 1,283 & 0.371 & High-flying tax lawter Cynthia Gibson Beerbower, 65, was found by her housekeeper at her apartment in Fulham, west London, in July this year . A high-flying tax lawyer who worked in the White House was found dead in the bath of her London home after taking an overdose while going through a 'painful' divorce. American-born Cynthia Gibson Beerbower, 65, was found by her housekeeper at her multi-million-pound apartment in Fulham, west London, on July 26 this year. The 'highly intelligent' former Wall Street lawyer had left three notes in her spare bedroom alongside a piece of card which read 'I can't do this anymore', an inquest has heard. The hearing was also told how the note included a request for Mrs Beerbower be buried in Cambridge, where she attended university. Mrs Beerbower had apparently been offered several ambassadorial roles by President Barack Obama in the weeks before her death and was planning a future with her new fiancé. But the hearing at Westminster Coroner's Court heard that Mrs Beerbower was struggling to cope with an ongoing divorce battle and feared her family 'wished her dead'. Her estranged husband also described how Mrs Beerbower had a history of mental illness and had threatened suicide several times. During the hearing, Coroner Fiona Wilcox told the court that a note left by the lawyer said: 'I can't do this any more, I hate you, I am sick'. The court also heard evidence from medical professionals who said Mrs Beerbower had started taking medication that she had not been prescribed. During the hearing, Dr Marios Pierides, a psychiatrist at the Nightingale Hospital in Marylebone, said Mrs Beerbower had visited him before her death. He described how she felt 'lonely' and 'isolated' while going through the process of divorcing her husband John Beerbower. He told the court: 'Cynthia was on various drugs at different time. 'When she came to me she was being treated with various medication, she was on anti-depressants and sleeping pills. 'She explained to me she was going through a very painful period of divorce. Also she was feeling quite lonely and isolated in her residence in Fulham. Dr Pierides added: 'She would say the family wished her dead. But she was, from my point of view, a life affirming person who enjoyed things until the awful circumstances around the divorce happened.' He continued: 'At the same time she was angry and lonely at her circumstances.' The court heard Mrs Beerbower had started taking Phentermine, an appetite suppressor, as well as Diazepam. Neither had been prescribed by the psychiatrist. Dr Pierides told the hearing that the drugs have the capability to 'alter' a person's mental state, adding: 'It could promote psychosis or mood swings. 'It can alter your mental state and you might not know what you're doing.' The American-born Cambridge graduate (left) had been offered several US ambassadorial roles by Barack Obama in the weeks before her death and was planning a future with her new partner, David Shaw (right) Mrs Beerbower's estranged husband, whom she married in 1971, told the court that Mrs Beerbower had threatened suicide in the past. He also said she had previously 'over-medicated' and had mental health problems dating back to the 1970s. Mr Beebower told the inquest: 'I think there were some issues that go all the way back to the middle of the 70s. It was between two and three decades that she had been seeing people. ' He said that she became suicidal on occasions, adding: 'She would become very upset or angry and would make these kinds of statements and, for example, lock herself in the bathroom. 'As far as I'm aware she never undertook any action on any of these threats.' The inquest heard Mrs Beerbower was struggling to cope with her divorce from husband John and feared her family 'wished her dead'. Pictured: Mr Beerbower and the couple's daughter Sarah outside Westminster Coroner's Court . Despite her divorce not being finalised, Mrs Beerbower become engaged to David Shaw in May this year, the inquest heard. Mr Shaw described her as a 'highly intelligent women' and said they had started to make plans for the future. He said: 'We were planning for the next 20 years. We were planning about what countries we wanted to go, because she was being offered ambassadorial roles by President Obama.' Mrs Beerbower - one of the few female partners in a Wall Street law firm - had joined the first Clinton administration in 1993. She served as International Tax Counsel under Secretary Lloyd Bentsen before becoming Deputy Assistant Secretary for Tax Policy under Secretary Robert Rubin. She left US Government service in late 1996. Mrs Beerbower retired from the financial business in 2005 and moved to France with her daughter Sarah. Mrs Beerbower was found dead by her housekeeper at her apartment in Fulham, west London (pictured) She and her husband acquired an olive farm in Chateauneuf de Grasse, in the south of France in 2006. From 2012 onward, Cynthia divided her time between homes in Toronto, New York City, the South of France and London, the inquest heard. Home Office pathologist Dr Jon Van Der Walt gave the cause of death as Oxidon, Phetermine and paracetemol toxicity. He said the drugs could have been enough to create a loss of consciousness, or fatal arrhythmia. The coroner said she could not give a verdict of suicide because the drugs Mrs Beerbower was taking could have affected her decision-making. Dr Wilcox said: 'At the time of death she was under the influence of Phentermine - it may well have affected her ability to have clear intentions, or even think clearly in any sensible way.' She instead recorded a narrative verdict, adding: 'She died as a result of poly-drug overdose while under the influence of Phentermine.' For confidential support call the Samaritans in the UK on 08457 90 90 90, visit a local Samaritans branch or click here for details. abound \\
\cline{1-9}
False & 42.5% & 12,929,214.5167 & \bfseries \color{red} 100.0000\% & 13.089 & 480.051 & 1,283 & 0.374 & High-flying tax lawter Cynthia Gibson Beerbower, 65, was found by her housekeeper at her apartment in Fulham, west London, in July this year . A high-flying tax lawyer who worked in the White House was found dead in the bath of her London home after taking an overdose while going through a 'painful' divorce. American-born Cynthia Gibson Beerbower, 65, was found by her housekeeper at her multi-million-pound apartment in Fulham, west London, on July 26 this year. The 'highly intelligent' former Wall Street lawyer had left three notes in her spare bedroom alongside a piece of card which read 'I can't do this anymore', an inquest has heard. The hearing was also told how the note included a request for Mrs Beerbower be buried in Cambridge, where she attended university. Mrs Beerbower had apparently been offered several ambassadorial roles by President Barack Obama in the weeks before her death and was planning a future with her new fiancé. But the hearing at Westminster Coroner's Court heard that Mrs Beerbower was struggling to cope with an ongoing divorce battle and feared her family 'wished her dead'. Her estranged husband also described how Mrs Beerbower had a history of mental illness and had threatened suicide several times. During the hearing, Coroner Fiona Wilcox told the court that a note left by the lawyer said: 'I can't do this any more, I hate you, I am sick'. The court also heard evidence from medical professionals who said Mrs Beerbower had started taking medication that she had not been prescribed. During the hearing, Dr Marios Pierides, a psychiatrist at the Nightingale Hospital in Marylebone, said Mrs Beerbower had visited him before her death. He described how she felt 'lonely' and 'isolated' while going through the process of divorcing her husband John Beerbower. He told the court: 'Cynthia was on various drugs at different time. 'When she came to me she was being treated with various medication, she was on anti-depressants and sleeping pills. 'She explained to me she was going through a very painful period of divorce. Also she was feeling quite lonely and isolated in her residence in Fulham. Dr Pierides added: 'She would say the family wished her dead. But she was, from my point of view, a life affirming person who enjoyed things until the awful circumstances around the divorce happened.' He continued: 'At the same time she was angry and lonely at her circumstances.' The court heard Mrs Beerbower had started taking Phentermine, an appetite suppressor, as well as Diazepam. Neither had been prescribed by the psychiatrist. Dr Pierides told the hearing that the drugs have the capability to 'alter' a person's mental state, adding: 'It could promote psychosis or mood swings. 'It can alter your mental state and you might not know what you're doing.' The American-born Cambridge graduate (left) had been offered several US ambassadorial roles by Barack Obama in the weeks before her death and was planning a future with her new partner, David Shaw (right) Mrs Beerbower's estranged husband, whom she married in 1971, told the court that Mrs Beerbower had threatened suicide in the past. He also said she had previously 'over-medicated' and had mental health problems dating back to the 1970s. Mr Beebower told the inquest: 'I think there were some issues that go all the way back to the middle of the 70s. It was between two and three decades that she had been seeing people. ' He said that she became suicidal on occasions, adding: 'She would become very upset or angry and would make these kinds of statements and, for example, lock herself in the bathroom. 'As far as I'm aware she never undertook any action on any of these threats.' The inquest heard Mrs Beerbower was struggling to cope with her divorce from husband John and feared her family 'wished her dead'. Pictured: Mr Beerbower and the couple's daughter Sarah outside Westminster Coroner's Court . Despite her divorce not being finalised, Mrs Beerbower become engaged to David Shaw in May this year, the inquest heard. Mr Shaw described her as a 'highly intelligent women' and said they had started to make plans for the future. He said: 'We were planning for the next 20 years. We were planning about what countries we wanted to go, because she was being offered ambassadorial roles by President Obama.' Mrs Beerbower - one of the few female partners in a Wall Street law firm - had joined the first Clinton administration in 1993. She served as International Tax Counsel under Secretary Lloyd Bentsen before becoming Deputy Assistant Secretary for Tax Policy under Secretary Robert Rubin. She left US Government service in late 1996. Mrs Beerbower retired from the financial business in 2005 and moved to France with her daughter Sarah. Mrs Beerbower was found dead by her housekeeper at her apartment in Fulham, west London (pictured) She and her husband acquired an olive farm in Chateauneuf de Grasse, in the south of France in 2006. From 2012 onward, Cynthia divided her time between homes in Toronto, New York City, the South of France and London, the inquest heard. Home Office pathologist Dr Jon Van Der Walt gave the cause of death as Oxidon, Phetermine and paracetemol toxicity. He said the drugs could have been enough to create a loss of consciousness, or fatal arrhythmia. The coroner said she could not give a verdict of suicide because the drugs Mrs Beerbower was taking could have affected her decision-making. Dr Wilcox said: 'At the time of death she was under the influence of Phentermine - it may well have affected her ability to have clear intentions, or even think clearly in any sensible way.' She instead recorded a narrative verdict, adding: 'She died as a result of poly-drug overdose while under the influence of Phentermine.' For confidential support call the Samaritans in the UK on 08457 90 90 90, visit a local Samaritans branch or click here for details. tournament \\
\cline{1-9}
False & 45.0% & 16,601,440.0572 & \bfseries \color{red} 100.0000\% & 12.079 & 477.168 & 1,283 & 0.372 & High-flying tax lawter Cynthia Gibson Beerbower, 65, was found by her housekeeper at her apartment in Fulham, west London, in July this year . A high-flying tax lawyer who worked in the White House was found dead in the bath of her London home after taking an overdose while going through a 'painful' divorce. American-born Cynthia Gibson Beerbower, 65, was found by her housekeeper at her multi-million-pound apartment in Fulham, west London, on July 26 this year. The 'highly intelligent' former Wall Street lawyer had left three notes in her spare bedroom alongside a piece of card which read 'I can't do this anymore', an inquest has heard. The hearing was also told how the note included a request for Mrs Beerbower be buried in Cambridge, where she attended university. Mrs Beerbower had apparently been offered several ambassadorial roles by President Barack Obama in the weeks before her death and was planning a future with her new fiancé. But the hearing at Westminster Coroner's Court heard that Mrs Beerbower was struggling to cope with an ongoing divorce battle and feared her family 'wished her dead'. Her estranged husband also described how Mrs Beerbower had a history of mental illness and had threatened suicide several times. During the hearing, Coroner Fiona Wilcox told the court that a note left by the lawyer said: 'I can't do this any more, I hate you, I am sick'. The court also heard evidence from medical professionals who said Mrs Beerbower had started taking medication that she had not been prescribed. During the hearing, Dr Marios Pierides, a psychiatrist at the Nightingale Hospital in Marylebone, said Mrs Beerbower had visited him before her death. He described how she felt 'lonely' and 'isolated' while going through the process of divorcing her husband John Beerbower. He told the court: 'Cynthia was on various drugs at different time. 'When she came to me she was being treated with various medication, she was on anti-depressants and sleeping pills. 'She explained to me she was going through a very painful period of divorce. Also she was feeling quite lonely and isolated in her residence in Fulham. Dr Pierides added: 'She would say the family wished her dead. But she was, from my point of view, a life affirming person who enjoyed things until the awful circumstances around the divorce happened.' He continued: 'At the same time she was angry and lonely at her circumstances.' The court heard Mrs Beerbower had started taking Phentermine, an appetite suppressor, as well as Diazepam. Neither had been prescribed by the psychiatrist. Dr Pierides told the hearing that the drugs have the capability to 'alter' a person's mental state, adding: 'It could promote psychosis or mood swings. 'It can alter your mental state and you might not know what you're doing.' The American-born Cambridge graduate (left) had been offered several US ambassadorial roles by Barack Obama in the weeks before her death and was planning a future with her new partner, David Shaw (right) Mrs Beerbower's estranged husband, whom she married in 1971, told the court that Mrs Beerbower had threatened suicide in the past. He also said she had previously 'over-medicated' and had mental health problems dating back to the 1970s. Mr Beebower told the inquest: 'I think there were some issues that go all the way back to the middle of the 70s. It was between two and three decades that she had been seeing people. ' He said that she became suicidal on occasions, adding: 'She would become very upset or angry and would make these kinds of statements and, for example, lock herself in the bathroom. 'As far as I'm aware she never undertook any action on any of these threats.' The inquest heard Mrs Beerbower was struggling to cope with her divorce from husband John and feared her family 'wished her dead'. Pictured: Mr Beerbower and the couple's daughter Sarah outside Westminster Coroner's Court . Despite her divorce not being finalised, Mrs Beerbower become engaged to David Shaw in May this year, the inquest heard. Mr Shaw described her as a 'highly intelligent women' and said they had started to make plans for the future. He said: 'We were planning for the next 20 years. We were planning about what countries we wanted to go, because she was being offered ambassadorial roles by President Obama.' Mrs Beerbower - one of the few female partners in a Wall Street law firm - had joined the first Clinton administration in 1993. She served as International Tax Counsel under Secretary Lloyd Bentsen before becoming Deputy Assistant Secretary for Tax Policy under Secretary Robert Rubin. She left US Government service in late 1996. Mrs Beerbower retired from the financial business in 2005 and moved to France with her daughter Sarah. Mrs Beerbower was found dead by her housekeeper at her apartment in Fulham, west London (pictured) She and her husband acquired an olive farm in Chateauneuf de Grasse, in the south of France in 2006. From 2012 onward, Cynthia divided her time between homes in Toronto, New York City, the South of France and London, the inquest heard. Home Office pathologist Dr Jon Van Der Walt gave the cause of death as Oxidon, Phetermine and paracetemol toxicity. He said the drugs could have been enough to create a loss of consciousness, or fatal arrhythmia. The coroner said she could not give a verdict of suicide because the drugs Mrs Beerbower was taking could have affected her decision-making. Dr Wilcox said: 'At the time of death she was under the influence of Phentermine - it may well have affected her ability to have clear intentions, or even think clearly in any sensible way.' She instead recorded a narrative verdict, adding: 'She died as a result of poly-drug overdose while under the influence of Phentermine.' For confidential support call the Samaritans in the UK on 08457 90 90 90, visit a local Samaritans branch or click here for details. Hybrid \\
\cline{1-9}
False & 47.5% & 8,886,110.5205 & \bfseries \color{red} 100.0000\% & 12.047 & 483.542 & 1,283 & 0.377 & High-flying tax lawter Cynthia Gibson Beerbower, 65, was found by her housekeeper at her apartment in Fulham, west London, in July this year . A high-flying tax lawyer who worked in the White House was found dead in the bath of her London home after taking an overdose while going through a 'painful' divorce. American-born Cynthia Gibson Beerbower, 65, was found by her housekeeper at her multi-million-pound apartment in Fulham, west London, on July 26 this year. The 'highly intelligent' former Wall Street lawyer had left three notes in her spare bedroom alongside a piece of card which read 'I can't do this anymore', an inquest has heard. The hearing was also told how the note included a request for Mrs Beerbower be buried in Cambridge, where she attended university. Mrs Beerbower had apparently been offered several ambassadorial roles by President Barack Obama in the weeks before her death and was planning a future with her new fiancé. But the hearing at Westminster Coroner's Court heard that Mrs Beerbower was struggling to cope with an ongoing divorce battle and feared her family 'wished her dead'. Her estranged husband also described how Mrs Beerbower had a history of mental illness and had threatened suicide several times. During the hearing, Coroner Fiona Wilcox told the court that a note left by the lawyer said: 'I can't do this any more, I hate you, I am sick'. The court also heard evidence from medical professionals who said Mrs Beerbower had started taking medication that she had not been prescribed. During the hearing, Dr Marios Pierides, a psychiatrist at the Nightingale Hospital in Marylebone, said Mrs Beerbower had visited him before her death. He described how she felt 'lonely' and 'isolated' while going through the process of divorcing her husband John Beerbower. He told the court: 'Cynthia was on various drugs at different time. 'When she came to me she was being treated with various medication, she was on anti-depressants and sleeping pills. 'She explained to me she was going through a very painful period of divorce. Also she was feeling quite lonely and isolated in her residence in Fulham. Dr Pierides added: 'She would say the family wished her dead. But she was, from my point of view, a life affirming person who enjoyed things until the awful circumstances around the divorce happened.' He continued: 'At the same time she was angry and lonely at her circumstances.' The court heard Mrs Beerbower had started taking Phentermine, an appetite suppressor, as well as Diazepam. Neither had been prescribed by the psychiatrist. Dr Pierides told the hearing that the drugs have the capability to 'alter' a person's mental state, adding: 'It could promote psychosis or mood swings. 'It can alter your mental state and you might not know what you're doing.' The American-born Cambridge graduate (left) had been offered several US ambassadorial roles by Barack Obama in the weeks before her death and was planning a future with her new partner, David Shaw (right) Mrs Beerbower's estranged husband, whom she married in 1971, told the court that Mrs Beerbower had threatened suicide in the past. He also said she had previously 'over-medicated' and had mental health problems dating back to the 1970s. Mr Beebower told the inquest: 'I think there were some issues that go all the way back to the middle of the 70s. It was between two and three decades that she had been seeing people. ' He said that she became suicidal on occasions, adding: 'She would become very upset or angry and would make these kinds of statements and, for example, lock herself in the bathroom. 'As far as I'm aware she never undertook any action on any of these threats.' The inquest heard Mrs Beerbower was struggling to cope with her divorce from husband John and feared her family 'wished her dead'. Pictured: Mr Beerbower and the couple's daughter Sarah outside Westminster Coroner's Court . Despite her divorce not being finalised, Mrs Beerbower become engaged to David Shaw in May this year, the inquest heard. Mr Shaw described her as a 'highly intelligent women' and said they had started to make plans for the future. He said: 'We were planning for the next 20 years. We were planning about what countries we wanted to go, because she was being offered ambassadorial roles by President Obama.' Mrs Beerbower - one of the few female partners in a Wall Street law firm - had joined the first Clinton administration in 1993. She served as International Tax Counsel under Secretary Lloyd Bentsen before becoming Deputy Assistant Secretary for Tax Policy under Secretary Robert Rubin. She left US Government service in late 1996. Mrs Beerbower retired from the financial business in 2005 and moved to France with her daughter Sarah. Mrs Beerbower was found dead by her housekeeper at her apartment in Fulham, west London (pictured) She and her husband acquired an olive farm in Chateauneuf de Grasse, in the south of France in 2006. From 2012 onward, Cynthia divided her time between homes in Toronto, New York City, the South of France and London, the inquest heard. Home Office pathologist Dr Jon Van Der Walt gave the cause of death as Oxidon, Phetermine and paracetemol toxicity. He said the drugs could have been enough to create a loss of consciousness, or fatal arrhythmia. The coroner said she could not give a verdict of suicide because the drugs Mrs Beerbower was taking could have affected her decision-making. Dr Wilcox said: 'At the time of death she was under the influence of Phentermine - it may well have affected her ability to have clear intentions, or even think clearly in any sensible way.' She instead recorded a narrative verdict, adding: 'She died as a result of poly-drug overdose while under the influence of Phentermine.' For confidential support call the Samaritans in the UK on 08457 90 90 90, visit a local Samaritans branch or click here for details. Santiago \\
\cline{1-9}
False & 50.0% & 7,366,844.3689 & \bfseries \color{red} 100.0000\% & 12.721 & 478.379 & 1,283 & 0.373 & High-flying tax lawter Cynthia Gibson Beerbower, 65, was found by her housekeeper at her apartment in Fulham, west London, in July this year . A high-flying tax lawyer who worked in the White House was found dead in the bath of her London home after taking an overdose while going through a 'painful' divorce. American-born Cynthia Gibson Beerbower, 65, was found by her housekeeper at her multi-million-pound apartment in Fulham, west London, on July 26 this year. The 'highly intelligent' former Wall Street lawyer had left three notes in her spare bedroom alongside a piece of card which read 'I can't do this anymore', an inquest has heard. The hearing was also told how the note included a request for Mrs Beerbower be buried in Cambridge, where she attended university. Mrs Beerbower had apparently been offered several ambassadorial roles by President Barack Obama in the weeks before her death and was planning a future with her new fiancé. But the hearing at Westminster Coroner's Court heard that Mrs Beerbower was struggling to cope with an ongoing divorce battle and feared her family 'wished her dead'. Her estranged husband also described how Mrs Beerbower had a history of mental illness and had threatened suicide several times. During the hearing, Coroner Fiona Wilcox told the court that a note left by the lawyer said: 'I can't do this any more, I hate you, I am sick'. The court also heard evidence from medical professionals who said Mrs Beerbower had started taking medication that she had not been prescribed. During the hearing, Dr Marios Pierides, a psychiatrist at the Nightingale Hospital in Marylebone, said Mrs Beerbower had visited him before her death. He described how she felt 'lonely' and 'isolated' while going through the process of divorcing her husband John Beerbower. He told the court: 'Cynthia was on various drugs at different time. 'When she came to me she was being treated with various medication, she was on anti-depressants and sleeping pills. 'She explained to me she was going through a very painful period of divorce. Also she was feeling quite lonely and isolated in her residence in Fulham. Dr Pierides added: 'She would say the family wished her dead. But she was, from my point of view, a life affirming person who enjoyed things until the awful circumstances around the divorce happened.' He continued: 'At the same time she was angry and lonely at her circumstances.' The court heard Mrs Beerbower had started taking Phentermine, an appetite suppressor, as well as Diazepam. Neither had been prescribed by the psychiatrist. Dr Pierides told the hearing that the drugs have the capability to 'alter' a person's mental state, adding: 'It could promote psychosis or mood swings. 'It can alter your mental state and you might not know what you're doing.' The American-born Cambridge graduate (left) had been offered several US ambassadorial roles by Barack Obama in the weeks before her death and was planning a future with her new partner, David Shaw (right) Mrs Beerbower's estranged husband, whom she married in 1971, told the court that Mrs Beerbower had threatened suicide in the past. He also said she had previously 'over-medicated' and had mental health problems dating back to the 1970s. Mr Beebower told the inquest: 'I think there were some issues that go all the way back to the middle of the 70s. It was between two and three decades that she had been seeing people. ' He said that she became suicidal on occasions, adding: 'She would become very upset or angry and would make these kinds of statements and, for example, lock herself in the bathroom. 'As far as I'm aware she never undertook any action on any of these threats.' The inquest heard Mrs Beerbower was struggling to cope with her divorce from husband John and feared her family 'wished her dead'. Pictured: Mr Beerbower and the couple's daughter Sarah outside Westminster Coroner's Court . Despite her divorce not being finalised, Mrs Beerbower become engaged to David Shaw in May this year, the inquest heard. Mr Shaw described her as a 'highly intelligent women' and said they had started to make plans for the future. He said: 'We were planning for the next 20 years. We were planning about what countries we wanted to go, because she was being offered ambassadorial roles by President Obama.' Mrs Beerbower - one of the few female partners in a Wall Street law firm - had joined the first Clinton administration in 1993. She served as International Tax Counsel under Secretary Lloyd Bentsen before becoming Deputy Assistant Secretary for Tax Policy under Secretary Robert Rubin. She left US Government service in late 1996. Mrs Beerbower retired from the financial business in 2005 and moved to France with her daughter Sarah. Mrs Beerbower was found dead by her housekeeper at her apartment in Fulham, west London (pictured) She and her husband acquired an olive farm in Chateauneuf de Grasse, in the south of France in 2006. From 2012 onward, Cynthia divided her time between homes in Toronto, New York City, the South of France and London, the inquest heard. Home Office pathologist Dr Jon Van Der Walt gave the cause of death as Oxidon, Phetermine and paracetemol toxicity. He said the drugs could have been enough to create a loss of consciousness, or fatal arrhythmia. The coroner said she could not give a verdict of suicide because the drugs Mrs Beerbower was taking could have affected her decision-making. Dr Wilcox said: 'At the time of death she was under the influence of Phentermine - it may well have affected her ability to have clear intentions, or even think clearly in any sensible way.' She instead recorded a narrative verdict, adding: 'She died as a result of poly-drug overdose while under the influence of Phentermine.' For confidential support call the Samaritans in the UK on 08457 90 90 90, visit a local Samaritans branch or click here for details. قاعدة \\
\cline{1-9}
\bottomrule
\end{tabular}
\end{table}
